% =============================================================================
% appendix.tex -- proofs and hypothesis lists
% Who Gets Caught: Disclosure Rules, Liquidity and the Undetected Blockholder
%
% Compile (from the repo root, after paper.tex and biber):
%   xelatex -interaction=nonstopmode appendix.tex
%
% Assembles the proof fragments under proofs/ by \input, in the order below:
% the partition subsection's standing conditions come first and every later
% fragment cites them by number; each fragment's cross-references stay inside
% the file that carries them.
% =============================================================================

\documentclass[11pt]{article}

\usepackage[margin=1in]{geometry}
\usepackage{amsmath,amssymb,amsthm}
\usepackage{setspace}
\usepackage{enumitem}
\usepackage[hidelinks]{hyperref}
\usepackage{etoolbox}

\setstretch{1.1}

\newtheorem{theorem}{Theorem}
\newtheorem{proposition}{Proposition}
\newtheorem{lemma}{Lemma}
\newtheorem{corollary}{Corollary}
\newtheorem{condition}{Condition}
\newtheorem{remark}{Remark}

\newif\ifunlabelled
\unlabelledfalse

\AtBeginEnvironment{theorem}{\ifunlabelled\else\par\noindent\textsc{Label: PROVED.}\par\nopagebreak\fi}
\AtBeginEnvironment{proposition}{\ifunlabelled\else\par\noindent\textsc{Label: PROVED.}\par\nopagebreak\fi}
\AtBeginEnvironment{corollary}{\ifunlabelled\else\par\noindent\textsc{Label: PROVED.}\par\nopagebreak\fi}
\AtBeginEnvironment{lemma}{\ifunlabelled\else\par\noindent\textsc{Label: PROVED.}\par\nopagebreak\fi}

\newcommand{\Eop}{\mathbb{E}}
\newcommand{\Prb}{\mathbb{P}}
\newcommand{\ind}{\mathbf{1}}

\title{\textbf{Who Gets Caught: Disclosure Rules, Liquidity and the Undetected Blockholder}\\[0.3em]
\large Appendix: statements, hypothesis lists, and proofs}
\author{Austin Li\\[0.25em]}
\date{This draft: 2 September 2026}

\begin{document}
\maketitle

\noindent This appendix carries every proof behind the paper's labelled
statements, and closes with two lemmas on the pricing root and the filing timing.
It starts with the disclosure partition and the fixed-policy results,
then gives the order-size-two pooled experiment, the threshold dial, and the
who-gets-caught corollary. The standing conditions (S1) to (S14) apply to
every statement that uses them. The corollary's hypotheses (C-1) to (C-3) cite
them by number.

\section{The disclosure partition and the fixed-policy results}
\label{sec:app-partition}

% ===========================================================================
% The flagged cell is free of the noise-trading intensity.
% ===========================================================================

\providecommand{\Eop}{\mathbb{E}}
\providecommand{\Prb}{\mathbb{P}}
\providecommand{\ind}{\mathbf{1}}
\providecommand{\ND}{\mathrm{ND}}
\providecommand{\Sfl}{\mathsf{S}_F}

\begin{lemma}[The flagged cell does not move with liquidity]
\label{lem:flagged-kappa-free}
Fix the cutoff policy and the execution policy, and let the disclosure rule $(\tau,T)$ be fixed.
Assume the primitive draws are mutually independent with strictly positive variances; the menu, the
image of the coarsening $\Gamma$, the noise support and the calendar are finite; every stake path
$s\mapsto B_j(s,d)$ is Borel, with Voice stakes weakly increasing in the signal and in the date; the
blockholder does not revise the stake path in response to realised order flow or prices; the legal-clock definitions
hold, including $D=1\Rightarrow a=1$ and the restriction that only Voice plans cross; the composed
terminal target $s\mapsto b^*_{j(s)}(s)$ is strictly increasing on the flagged signal region, almost
surely on the flagged set; the inner pricing fixed point is unique; the flagged weight satisfies
$\Omega>0$; and the bidder's entry rule has the two properties
\begin{equation}
\label{eq:entry-props}
\Prb(\mathsf B=1\mid\mathcal I_H)=p(\mathcal I_H)
\qquad\text{and}\qquad
\mathsf B \perp v \mid \mathcal I_H .
\end{equation}
Then, on the flagged cell $\mathcal C_F$, the flagged tuple $\Sfl=(B^F,Q^F,a{=}1)$ makes the
pre-filing pooled history conditionally independent of $(v,s,\xi)$, and the flagged posterior, the
flagged price, the bidder's entry probability and the flagged cell average $M_F$ do not depend on
$\kappa$. In addition $\Omega$ does not depend on $\kappa$, so $\partial_\kappa\Omega=0$ at fixed
policies.
\end{lemma}

\begin{proof}
Write $\Xi:=(v,s,\xi)$ for the triple whose conditional independence is at issue, $z_{0:H}$ for the
noise vector, $\mathcal H_{f^-}^P$ for the pre-filing pooled history on a flagged path, and
$\mathcal I_H$ for the public information set at the control node, which the model sets equal to
$\Sfl$ on the flagged cell. The functions $u_1,u_2,u_3,u_5$ below are bounded and measurable and are
local to this proof.

\emph{Step 1 (where $\kappa$ lives).} At fixed cutoff and execution policies, and with no revision of
the path within the window, every policy object is a function of $(s;\tau,T)$ alone: the selected
plan $j=j(s)$, the stake path $B_j(s,\cdot)$, the order marks $q_{jd}(s)$, the terminal target
$b_j^*(s)$, the crossing date $c_j(s)$, the filing date $f_j(s)$, the disclosure indicator $D_j(s)$,
the stake at filing $B_j^F(s)$ and the flagged order $Q_j^F(s)$. Absence of feedback removes any
dependence on realised order flow or prices, and fixed policies remove any dependence on $\kappa$
through re-optimised cutoffs. Among the primitives, $\kappa$ appears in exactly one place, the law of
the ternary noise mark; the laws of $v$, $\varepsilon$ and $\xi$ and the remaining constants carry no
$\kappa$. Write $Q_\kappa$ for the law of $z_{0:H}$. At fixed policies, then, $\kappa$ enters the
model only through $Q_\kappa$.

\emph{Step 2 (the flagged set and its mass are free of $\kappa$).} The disclosure indicator
$D=\ind\{a=1,\ c(\tau)+T\le H\}$ is measurable, and by Step~1 it is a function of the signal alone.
Hence the flagged cell $\mathcal C_F=\{D=1\}$ lies in $\sigma(s)$ and is one fixed Borel subset of
the signal line, the same subset at every $\kappa$. The signal $s=v+\varepsilon$ is
$N(\mu_v,\sigma_v^2+\sigma_\varepsilon^2)$, a law free of $\kappa$ because the noise mark enters
neither $v$ nor $\varepsilon$. Therefore $\Omega=\Prb(s\in\mathcal C_F)$ is the same number at every
$\kappa$, which is the last conclusion of the statement.

\emph{Step 3 (the flagged tuple pins the signal, with a measurable inverse).} On $\mathcal C_F$ the
engagement flag is $a=1$, since disclosure implies engagement and only Voice plans cross, so the
third coordinate of $\Sfl$ is uninformative there. By the definition $Q_j^F=b_j^*(s)-B_j^F$,
\begin{equation}
\label{eq:flagged-target}
B^F_{j(s)}(s) \;+\; Q^F_{j(s)}(s) \;=\; b^*_{j(s)}(s).
\end{equation}
The composed terminal target on the right of \eqref{eq:flagged-target} is strictly increasing in the
signal on the flagged region and is therefore injective there. That monotonicity is consumed almost
surely on the flagged set, which is the only coherent reading when $B^F$ is continuum-valued: no
individual flagged tuple then carries positive probability, so injectivity holds almost surely on a
set of positive probability rather than on a set of atoms. Consequently the first two coordinates of
$\Sfl$ determine $b^*_{j(s)}(s)$, which determines $s$, which at fixed policies determines $j=j(s)$.
Write $\iota_F$ for the resulting inverse. It is strictly increasing on the image of the composed
target, and a monotone real function on a Borel set is Borel, so $\iota_F$ is measurable; no
measurable-selection theorem is needed. Since $\Omega>0$, the set on which this holds carries
strictly positive probability and the statement is not vacuous. Restricted to $\mathcal C_F$,
therefore, $\sigma(\Sfl)=\sigma(s)$ up to null sets. Strict monotonicity is what the step consumes:
if two signals produced the same flagged tuple while carrying different pooled paths,
\eqref{eq:flagged-target} would fail to be injective and the tuple would not pin the pooled path.

\emph{Step 4 (the pre-filing pooled history is a noise transform indexed by plan and signal).}
Define, for a plan $j$ and a signal $s$,
\begin{equation}
\label{eq:flagged-upsilon}
\Upsilon_{j,s}(z_{0:H}) \;:=\;
  \Bigl( \bigl(q_{j0}(s)+z_0,\dots,q_{j,f_j(s)-1}(s)+z_{f_j(s)-1}\bigr);\ \text{flag not landed by }
  f_j(s)-1 \Bigr),
\end{equation}
so that $\mathcal H_{f^-}^P=\Upsilon_{j(s),s}(z_{0:H})$ pointwise on $\mathcal C_F$. The indexing by
$(j,s)$ alone is exactly the absence of feedback: with re-optimisation inside the window the marks
would depend on realised flow, hence on past noise, and no such indexing would exist. The map
\eqref{eq:flagged-upsilon} is jointly measurable in $(s,z_{0:H})$. The filing date takes finitely
many values because the calendar is finite, and on each of its finitely many measurable level sets
the map is $(s,z)\mapsto(q_{j(s)d}(s)+z_d)_{d<t}$, measurable because
$q_{jd}=\Gamma\bigl(B_j(\cdot,d)-B_j(\cdot,d-1)\bigr)$ composes a Borel function, Voice stakes being
monotone in the signal and only Voice plans reaching $\mathcal C_F$, with the coarsening $\Gamma$,
itself Borel because it is ordered in the increment. This is unaffected by the size of the
blockholder's per-round order, which fixes the value the mark takes on a building round but not its
measurability in the signal. Patching finitely many measurable restrictions gives a measurable map.
The second component of \eqref{eq:flagged-upsilon} is a deterministic function of $s$ by Step~1 and
carries no information beyond $s$.

\emph{Step 5 (observed prices enlarge nothing).} Each pooled price is the unique fixed point of the
pricing map at its pooled history and is therefore a function of that history, and likewise the
flagged price is a function of the flagged information set. Adjoining the value of a function of a
generator of a $\sigma$-algebra to that $\sigma$-algebra does not enlarge it, so prices may be
ignored when the information content of $\mathcal I_H$ is computed.

\emph{Step 6 (an information sandwich).} On $\mathcal C_F$,
\begin{equation}
\label{eq:flagged-sandwich}
\sigma(s) \cap \mathcal C_F \;\subseteq\; \mathcal I_H \;\subseteq\;
  \sigma(s,z_{0:H}) \cap \mathcal C_F .
\end{equation}
For the left inclusion: the filing reveals $F=(B^F,a{=}1)$ and the flagged order is submitted and
priced, so $\Sfl$ belongs to $\mathcal I_H$; by Step~3 the tuple recovers $s$; and
$\mathcal C_F\in\sigma(s)$ by Step~2, so restricting to the flagged cell is itself an operation
within $\sigma(s)$. For the right inclusion: every component of $\mathcal I_H$, namely pooled order
flow, flag status, filing content, the flagged order and prices, is a measurable function of
$(s,z_{0:H})$ by Steps~1, 4 and~5. Everything below uses only \eqref{eq:flagged-sandwich}, so the
conclusion is robust to the exact content assigned to the control-node information set, provided it
satisfies the sandwich.

\emph{Step 7 (a freezing lemma).} Let $\mathcal G$ be a $\sigma$-algebra with $z_{0:H}$ independent
of $\mathcal G$, let $S$ be $\mathcal G$-measurable, let $w$ be bounded and jointly measurable, and
put $\bar w(x):=\int w(x,\zeta)\,Q_\kappa(\mathrm d\zeta)$. Then
$\Eop[w(S,z_{0:H})\mid\mathcal G]=\bar w(S)$. For a product form $w(x,\zeta)=u_1(x)u_2(\zeta)$ this
is $\Eop[u_1(S)u_2(z_{0:H})\mid\mathcal G]=u_1(S)\Eop[u_2(z_{0:H})]=\bar w(S)$, pulling out the
$\mathcal G$-measurable factor and using independence. The product forms generate the product
$\sigma$-algebra and are closed under multiplication; both sides of the asserted equality are linear
in $w$ and stable under bounded monotone limits, so the monotone class theorem extends the equality
to every bounded jointly measurable $w$.

\emph{Step 8 (the conditional independence).} By Step~3, conditioning on $\sigma(\Sfl)$ on the
flagged set is conditioning on $\sigma(s)$, so what has to be shown is
$\mathcal H_{f^-}^P \perp \Xi \mid s$ on $\mathcal C_F$. Given $s$, the middle coordinate of $\Xi$ is
degenerate, so it suffices to show, for bounded measurable $u_1$ on the pre-filing history space and
bounded measurable $u_2$ on $\mathbb R^2$,
\begin{equation}
\label{eq:flagged-factor}
\Eop\bigl[u_1(\mathcal H_{f^-}^P)\,u_2(v,\xi) \mid s\bigr]
 \;=\; \Eop\bigl[u_1(\mathcal H_{f^-}^P) \mid s\bigr]\cdot\Eop\bigl[u_2(v,\xi) \mid s\bigr]
 \qquad \text{a.s.\ on } \mathcal C_F,
\end{equation}
which is the standard characterisation of conditional independence given a $\sigma$-algebra. Put
$w(x,\zeta):=u_1\bigl(\Upsilon_{j(x),x}(\zeta)\bigr)$, bounded and jointly measurable by Step~4, so
that $u_1(\mathcal H_{f^-}^P)=w(s,z_{0:H})$. The noise vector is independent of
$(v,\varepsilon,\xi)$ and hence of $\Xi$, which is a measurable function of that triple. Apply
Step~7 with $\mathcal G=\sigma(v,s,\xi)$ and $S=s$:
$\Eop[w(s,z_{0:H})u_2(v,\xi)\mid v,s,\xi]=u_2(v,\xi)\,\bar w(s)$. Taking $\Eop[\,\cdot\mid s]$ of
both sides and pulling out the $\sigma(s)$-measurable factor $\bar w(s)$ gives
$\Eop[u_1u_2\mid s]=\bar w(s)\,\Eop[u_2\mid s]$; setting $u_2\equiv1$ in the same display gives
$\Eop[u_1\mid s]=\bar w(s)$, and substituting yields \eqref{eq:flagged-factor}. Unconditionally the
pre-filing pooled history depends strongly on the signal through the order marks, and the
conditioning event removes that dependence by pinning the signal, leaving the noise as the only
remaining randomness in \eqref{eq:flagged-upsilon}.

\emph{Step 9 (the flagged posterior).} Let $u_3$ be bounded, $\mathcal I_H$-measurable and supported
on $\mathcal C_F$. By the right inclusion of \eqref{eq:flagged-sandwich} and the Doob--Dynkin lemma,
$u_3=u_5(s,z_{0:H})$ for some bounded jointly measurable $u_5$. Running the computation of Step~8
with $u_5$ in place of $w$ gives $\Eop[u_2u_3]=\Eop\bigl[\bar u_5(s)\,\Eop[u_2\mid s]\bigr]$, while
the tower property together with $\Eop[u_3\mid s]=\bar u_5(s)$ from Step~7 gives
$\Eop\bigl[u_3\,\Eop[u_2\mid s]\bigr]=\Eop\bigl[\bar u_5(s)\,\Eop[u_2\mid s]\bigr]$. Hence
$\Eop[u_2u_3]=\Eop\bigl[\Eop[u_2\mid s]\,u_3\bigr]$ for every such $u_3$; since $\Eop[u_2\mid s]$ is
$\sigma(s)$-measurable and the left inclusion of \eqref{eq:flagged-sandwich} places
$\sigma(s)\cap\mathcal C_F$ inside $\mathcal I_H$, it is a version of $\Eop[u_2\mid\mathcal I_H]$ on
the flagged set:
\begin{equation}
\label{eq:flagged-posterior}
\Eop[u_2(v,\xi) \mid \mathcal I_H] \;=\; \Eop[u_2(v,\xi) \mid s]
 \qquad \text{a.s.\ on } \mathcal C_F .
\end{equation}
The extension from bounded to integrable $u_2$, needed at Step~11 for $u_2(v,\xi)=v$, follows by
applying \eqref{eq:flagged-posterior} to the truncations $\max(-n,\min(n,u_2))$ and passing to the
limit under conditional dominated convergence, the dominating variable being $|u_2|$, integrable
because $v$ is Gaussian. Separately, the engagement posterior satisfies $\pi(\mathcal I_H)=1$ on
$\mathcal C_F$, because the flagged cell lies inside $\{a=1\}$, and because the cell is itself
$\mathcal I_H$-measurable, being $\{D=1\}$ with $D$ part of the flag status carried in
$\mathcal I_H$.

\emph{Step 10 (the posterior does not move with $\kappa$).} By \eqref{eq:flagged-posterior} the
conditional law of $(v,\xi)$ given $\mathcal I_H$ on the flagged set equals its conditional law given
$s$. That law is $v\mid s \sim N\bigl(\mu_v+\beta(s-\mu_v),(1-\beta)\sigma_v^2\bigr)$ with
$\beta=\sigma_v^2/(\sigma_v^2+\sigma_\varepsilon^2)$, and $\xi\mid s \sim N(0,\sigma_\xi^2)$
independent of $v$. Neither depends on $\kappa$, which by Step~1 lives only in $Q_\kappa$ and enters
no coordinate of $\Xi$. With $\pi=1$ from Step~9, the flagged posterior over $(v,a,\xi)$ is a
function of the signal that does not move with $\kappa$, hence by Step~3 a function of $\Sfl$ that
does not move with $\kappa$.

\emph{Step 11 (the flagged price).} The flagged price solves the inner fixed point
$P=\Eop[Y\mid\mathcal I_H]$. On the flagged cell $a=1$ and $\pi=1$, so
$m_0+\Delta_m=m_1$. Only the two entry properties \eqref{eq:entry-props} are used. The entry rule of
the model has both, since the bidder's shock satisfies $\xi\mid\mathcal I_H \sim N(0,\sigma_\xi^2)$
independently of $v$ by Step~10 and the entry indicator is a function of $\xi$ and of
$\mathcal I_H$-measurable quantities; any other rule with these two properties works as well. Under
them,
\begin{equation}
\label{eq:flagged-innerFP}
P \;=\; \bigl(1-p(P)\bigr)\Bigl(\mu_v+\beta(s-\mu_v)+\Delta_V\Bigr) \;+\; p(P)\bigl(P+m_1\bigr),
\qquad
p(P) \;=\; 1-\Phi\!\left(\frac{P+K+m_1-\bar S}{\sigma_\xi}\right),
\end{equation}
where $\Eop[v\mid\mathcal I_H]=\Eop[v\mid s]=\mu_v+\beta(s-\mu_v)$ by
\eqref{eq:flagged-posterior}. The data of \eqref{eq:flagged-innerFP} are the single number
$\Eop[v\mid s]$ and the primitive parameters, none of which depends on $\kappa$ by Step~10.
Uniqueness of the inner fixed point makes the solution a function of the data alone, so equal data
give equal solutions. Uniqueness is load-bearing rather than decorative: without it an extraneous
selection among fixed points could vary with $\kappa$ even though the map does not. Hence the flagged
price is a function of the signal, and by Step~3 of $\Sfl$, that does not move with $\kappa$.

\emph{Step 12 (the entry probability).} With $\pi=1$ from Step~9, the flagged entry probability is
$p=1-\Phi\bigl((P^F+K+m_1-\bar S)/\sigma_\xi\bigr)$ on the flagged cell. Every argument is free of
$\kappa$, by Step~11 for the price and by the primitive table for the constants. Write $p_F(s)$ for
it. It is measurable as a composition of the monotone map $s\mapsto\Eop[v\mid s]$, the fixed-point
solution and the normal distribution function, and it takes values in $(0,1)$.

\emph{Step 13 (the flagged cell average).} The kernel of the premium is $h=\pi p$, so $h=p_F(s)$ on
the flagged cell by Steps~9 and~12. Since $\Omega>0$, the elementary definition of conditioning on
a positive-probability event gives
\begin{equation}
\label{eq:flagged-MF}
M_F \;=\; \Delta_m\,\Eop[h \mid D=1]
     \;=\; \frac{\Delta_m\,\Eop\bigl[p_F(s)\,\ind\{s \in \mathcal C_F\}\bigr]}{\Omega} .
\end{equation}
The numerator of \eqref{eq:flagged-MF} integrates a bounded measurable function free of $\kappa$
(Step~12) over a Borel set free of $\kappa$ (Step~2) against the law of the signal, itself free of
$\kappa$ (Step~2), and the denominator is free of $\kappa$ and strictly positive. Hence $M_F$ does
not move with $\kappa$, which completes the four invariance conclusions.
\end{proof}

Two boundaries of the result are part of it. Nothing above asserts that the pre-filing pooled price
does not move with $\kappa$: that price conditions on order flow whose law is $Q_\kappa$, and it
moves with $\kappa$ in general. With $J=P^F-P_{\ND}$, where $P_{\ND}$ is the last pre-filing pooled
price at the same realised order flow, and with $\partial_\kappa P^F=0$ on the flagged set, wherever
the derivative exists $\partial_\kappa J=-\partial_\kappa P_{f^-}^P$. The result therefore makes no
claim about the filing-day jump. And the result is a fixed-policy statement: if cutoffs or execution
paths move with $\kappa$, Step~1 is unavailable and the conclusions do not follow.

%==============================================================================
% The partition and the factorisation  S = (1 - Omega) S_P
%==============================================================================

\providecommand{\Prb}{\mathbb{P}}
\providecommand{\Eop}{\mathbb{E}}
\providecommand{\ind}{\mathbf{1}}
\providecommand{\Dact}{\Delta_{\mathrm{act}}}

\subsection{The partition and the factorisation}
\label{sec:partition-factorisation}

The disclosure rule $(\tau,T)$ acts on the model first as a partition of histories and only then as
a statement about prices. This subsection separates the two. The partition is combinatorial and
needs no restriction on cell mass; the two-cell decomposition of the premium needs an interior
partition; the factorisation of the liquidity sensitivity needs, in addition, that the flagged
endpoint and the flagged weight are both free of $\kappa$.

\subsubsection*{Standing conditions}

The following are in force throughout the subsection.

\begin{enumerate}[label=(S\arabic*),leftmargin=3.0em,itemsep=2pt]
\item \label{p:space} \emph{One probability space.} The primitive vector, the value $v$, the signal
  noise $\varepsilon$, the bidder draw $\xi$ and the noise-trader marks $z_{0:H}$, has a joint law
  on a finite product of Polish spaces, with the noise marks valued in the finite set
  $\{-\bar z,0,+\bar z\}$. The signal is $s = v + \varepsilon$.
\item \label{p:menu} \emph{A finite menu and a finite calendar.} The plan menu $\mathcal J$ is
  finite, the horizon $H$ is finite, and the window margin satisfies $T \in \{1,\dots,H\}$.
\item \label{p:cutoff} \emph{A cutoff selection map.} The signal is carried into a plan by a Borel step
  function $j(\cdot)$ with breakpoints $k_1 \le \dots \le k_{J-1}$.
\item \label{p:paths} \emph{Monotone Voice paths and a clean start.} Each plan $j$ carries a Borel stake
  path $B_j(s,\cdot)$ with $B_j(s,-1) = b_0 < \tau$; for a Voice plan the path is weakly increasing
  in the signal and weakly increasing in the calendar date.
\item \label{p:clock} \emph{Legal-clock discipline.} Only Voice plans cross the threshold; the
  crossing date $c_j(s)$ is the first date at which the path reaches $\tau$, or $+\infty$ if the
  path never reaches it; the filing lands exactly at $f_j = c_j + T$ and only through the disclosure
  node; the filing reports the stake truthfully.
\item \label{p:public} \emph{The flag is public.} Every control-node history carries the public
  coordinate ``the filing has landed by date $d$'', so the control-node information set
  $\mathcal I_H$ contains that coordinate at $d = H$ on both cells. On the pooled cell
  $\mathcal I_H$ also contains the pooled history up to the control node.
\item \label{p:nofeedback} \emph{No-feedback timing.} The executed path, the order marks, the
  terminal target, the crossing date, the filing date, the filing stake and the flagged order are
  Borel functions of the plan and the signal alone, with the filing stake and flagged order defined
  on the flagged cell. Neither realised order flow nor a realised price enters any of them.
\item \label{p:kernel} \emph{A bounded, pinned kernel.} The engagement-premium kernel is
  $h(\mathcal I) = \pi(\mathcal I)\,p(\mathcal I)$, where the engagement posterior
  $\pi(\mathcal I) = \Prb(a = 1 \mid \mathcal I)$ takes values in $[0,1]$ and the bidder-entry
  probability $p(\mathcal I)$ takes values in $(0,1)$ and is a continuous function of the posterior
  and the price. The pricing rule pins one version of $\Eop[Y \mid \mathcal I]$ rather than an
  equivalence class.
\item \label{p:wedge} \emph{A finite wedge.} The premium wedge $\Delta_m$ is a finite, strictly
  positive constant, and $\Dact = \Delta_m\,\Eop[h(\mathcal I_H)]$.
\item \label{p:kappa} \emph{Liquidity enters in one place.} The intensity $\kappa$ appears in the
  primitives only in the law of the ternary noise mark. The laws of $v$, $\varepsilon$ and $\xi$
  and the remaining constants carry no $\kappa$.
\item \label{p:fixed} \emph{Fixed policies.} The plan menu, the execution policies and the cutoff
  vector $k$ are held fixed in $\kappa$ and held fixed across any two rules compared.
For the remaining conditions, let $\theta=(j(s),s)$ denote the blockholder's type,
$m_d(\theta)=q_{j(s)d}(s)$ its order mark at date $d$, and $r$ an arbitrary disclosure rule under
comparison. Write $\mathsf S_{F,r}=(B_r^F,Q_r^F,a{=}1)$ for that rule's flagged tuple.
\item \label{p:noise} \emph{A common noise channel.} The vector $z_{0:H}$ is independent of
  $(v,\varepsilon,\xi)$ and has the common rule-invariant law $Q_\kappa$. Pooled flow is
  $X_d=m_d(\theta)+z_d$. In the order-size-two model,
  $m_d(\theta)\in\{0,2\bar z\}$, and the noise marks are independent across dates and have
  probabilities $(\kappa/2,1-\kappa,\kappa/2)$ on
  $\{-\bar z,0,+\bar z\}$.
\item \label{p:control} \emph{The control-node convention.} On the flagged cell the information set
  is exactly $\sigma(\mathsf S_{F,r})$, up to pinned price coordinates. On the pooled cell it is
  exactly the displayed public pooled history, again up to pinned prices. At each depth $d$, the
  same convention applies to the cells $\mathcal C_{F,r}^d$ and $\mathcal C_{P,r}^d$ defined by
  whether the filing has landed by $d$.
\item \label{p:decoder} \emph{Identified flagged types.} On rule $r$'s flagged signal region, the
  composed terminal target $g(s)=b^*_{j(s)}(s)$ is strictly increasing and
  $B^F_r+Q^F_r=g(s)$. Pointwise strict increase is imposed for pointwise experiment comparisons;
  for posterior conclusions, strict increase outside a prior-null set is enough.
\end{enumerate}

\subsubsection*{The partition}

\begin{lemma}[Disclosure partition]
\label{lem:partition}
Under \ref{p:space}--\ref{p:public} and \ref{p:control}, the disclosure indicator
$D = \ind\{a = 1,\ c(\tau) + T \le H\}$ is a measurable function of the control-node history, and
the flagged cell $\mathcal C_F = \{D = 1\}$ and the pooled cell $\mathcal C_P = \{D = 0\}$ are
disjoint and cover the space. Consequently the flagged weight $\Omega = \Prb(D = 1)$ is defined and
$\Prb(D = 0) = 1 - \Omega$. No restriction on the mass of either cell is used.
\end{lemma}

\begin{proof}
\emph{Step 1 (the selection map is Borel).} By \ref{p:space} the signal $s = v + \varepsilon$ is a
sum of two coordinates of a Borel vector and is therefore Borel. Condition \ref{p:cutoff} makes
$j(\cdot)$ Borel. Its step form has finitely many breakpoints by \ref{p:menu}. The set
$\{s : j(s) = j\}$ is Borel for each
$j \in \mathcal J$.

\emph{Step 2 (the engagement flag is measurable).} The realised flag is $a = a_{j(s)}$, the
composition of the Borel map of Step~1 with the map $j \mapsto a_j$ defined on the finite set
$\mathcal J$. A map on a finite set carrying its discrete $\sigma$-algebra is measurable, so
$\{a = 1\}$ is Borel and lies in $\sigma(s)$.

\emph{Step 3 (the crossing date is attained, unique and measurable).} Fix a Voice plan $j$. By
\ref{p:paths} the map $s \mapsto B_j(s,d)$ is weakly increasing, and the preimage of
$[\tau,\infty)$ under a weakly monotone real function is an interval, so
$\{s : B_j(s,d) \ge \tau\}$ is Borel at each date $d$. The calendar $\{0,\dots,H\}$ is finite by
\ref{p:menu}, so
\[
  \{c_j \le d\} \;=\; \bigcup_{d' \le d}\{s : B_j(s,d') \ge \tau\}
\]
is a finite union of Borel sets and $c_j$ is a Borel map into $\{0,\dots,H\} \cup \{+\infty\}$. Its
value is unique: a nonempty subset of the finite well-ordered set $\{0,\dots,H\}$ has exactly one
minimum, and by \ref{p:clock} the crossing date is that minimum; if the subset is empty then
$c_j = +\infty$.

\emph{Step 4 (the disclosure indicator is measurable).} By \ref{p:clock} a filing requires
$a = 1$, and $f_j = c_j + T \le H$ holds exactly when $c_j \le H - T$, the convention
$(+\infty) + T = +\infty > H$ being available because $H$ is finite by \ref{p:menu}. Hence
\begin{equation}
\label{eq:cell-union}
  \{D = 1\} \;=\; \bigcup_{j \in \mathcal J\,:\,a_j = 1}
    \Bigl( \{s : j(s) = j\} \;\cap\; \{s : c_j(s) \le H - T\} \Bigr).
\end{equation}
Each set on the right is Borel by Steps~1 and~3, and the union is finite by \ref{p:menu}, so
$\{D = 1\}$ is Borel and $D$ is a measurable $\{0,1\}$-valued map. Record one fact for later: at a
fixed cutoff policy every ingredient of \eqref{eq:cell-union} is a function of $s$ alone, so
$\{D = 1\} \in \sigma(s)$. Record a second: regularity of $s \mapsto B_j(s,d)$ was needed only for
Voice plans, the conjunct $a = 1$ having removed the others from the union.

\emph{Step 5 (exactly one cell, at the level of states).} By Step~4, $D$ is the indicator of a Borel
event, so at every sample point it takes exactly one of the values $0$ and $1$. Therefore
$\mathcal C_F \cap \mathcal C_P = \varnothing$ and $\mathcal C_F \cup \mathcal C_P$ is the whole
space.

\emph{Step 6 (the cell is a function of the observed history).} Step~5 partitions the space of
states; the claim is about control-node histories. By \ref{p:public} every control-node history carries the
public coordinate ``the filing has landed by date $d$'', and
$\mathcal I_H$ contains that coordinate at $d = H$ on both cells. Under \ref{p:control}, flagged
tuples and pooled histories occupy disjoint tagged output spaces. By \ref{p:clock} the filing lands exactly at
$c + T$ and only through the disclosure node, so the coordinate ``the filing has landed by $H$''
equals $\ind\{f \le H\}$; combining this with $D = \ind\{a = 1,\ c < \infty,\ f \le H\}$ and with the
restriction that only Voice plans cross, that coordinate equals $D$. Hence $D$ is measurable with
respect to the control-node history and the map from a history to its cell is single-valued.
When both cells have positive probability, the map is onto
$\{\mathcal C_F, \mathcal C_P\}$. When $\Omega \in \{0,1\}$, one cell is null and the map
hits only the other almost surely, as Step~7 records. The public flag coordinate is what carries
this step: without it
$D$ would still be a measurable function of the state by Step~4, while the cell assignment would
stop being a function of what is observed at the control node.

\emph{Step 7 (the weight).} By Steps~5 and~6 the two cells are measurable, disjoint and exhaustive,
so $\Omega = \Prb(D = 1)$ is defined and $\Prb(D = 0) = 1 - \Omega$. Steps~1 to~7 used no
restriction on the value of $\Omega$, which may be $0$ or $1$.
\end{proof}

\subsubsection*{The two-cell decomposition}

\begin{lemma}[Two-cell decomposition of the premium]
\label{lem:two-cell}
Under \ref{p:space}--\ref{p:wedge}, with the partition of Lemma~\ref{lem:partition}, and whenever
$0 < \Omega < 1$,
\begin{equation}
\label{eq:two-cell}
  \Dact \;=\; \Omega\,M_F \;+\; (1 - \Omega)\,M_P,
  \qquad
  M_F = \Delta_m\Eop[h \mid D = 1],
  \quad
  M_P = \Delta_m\Eop[h \mid D = 0].
\end{equation}
At $\Omega = 1$ the identity reads $\Dact = M_F$ and at $\Omega = 0$ it reads $\Dact = M_P$; in each
of those two cases the average over the null cell is undefined rather than zero, and it is not
determined by the law. That $M_P$ is defined as a number, rather than an equivalence class, rests
on \ref{p:kernel} pinning a version of the price; Borel regularity of the stake paths in
\ref{p:paths} does not by itself select a version of a conditional expectation.
\end{lemma}

\begin{proof}
\emph{Step 1 (the kernel is bounded).} By \ref{p:kernel}, $\pi \in [0,1]$ and $p \in (0,1)$, so
their product satisfies $0 \le h \le 1$ pointwise.

\emph{Step 2 (the kernel is a random variable).} The posterior
$\pi(\mathcal I_H) = \Prb(a = 1 \mid \mathcal I_H)$ is a version of a conditional probability and is
$\mathcal I_H$-measurable, and so is the price $\Eop[Y \mid \mathcal I_H]$, whose version is pinned
by \ref{p:kernel}. That is what makes $h(\mathcal I_H)$ one well-defined random variable rather than
an equivalence class from which the decomposition would have to choose. The entry probability is the
composition of the posterior and the price with a continuous map, hence $\mathcal I_H$-measurable,
and a product of measurable functions is measurable. With Step~1 the kernel is integrable.

\emph{Step 3 (the pointwise split).} By Lemma~\ref{lem:partition} exactly one of $\ind\{D = 1\}$ and
$\ind\{D = 0\}$ equals one at every sample point and the other equals zero, so
$h = h\,\ind\{D = 1\} + h\,\ind\{D = 0\}$ pointwise.

\emph{Step 4 (integration).} Each summand in Step~3 is bounded in absolute value by $1$ by Step~1,
hence integrable on the single probability space of \ref{p:space}. Linearity of the integral gives
$\Eop[h] = \Eop[h\,\ind\{D = 1\}] + \Eop[h\,\ind\{D = 0\}]$.

\emph{Step 5 (the interior case).} Suppose $0 < \Omega < 1$. Both conditioning events have strictly
positive probability, so the elementary conditional expectations
$\Eop[h \mid D = 1] := \Eop[h\,\ind\{D = 1\}]/\Omega$ and
$\Eop[h \mid D = 0] := \Eop[h\,\ind\{D = 0\}]/(1 - \Omega)$ are defined, and Step~1 bounds both by
$1$. This is the defining property of conditioning on an event of positive probability. Substituting
into Step~4 gives $\Eop[h] = \Omega\,\Eop[h \mid D = 1] + (1 - \Omega)\,\Eop[h \mid D = 0]$.

\emph{Step 6 (scaling to premia).} By \ref{p:wedge} the wedge $\Delta_m$ is finite and strictly
positive. Multiplying Step~5 by $\Delta_m$ and reading off $\Dact$, $M_F$ and $M_P$ gives
\eqref{eq:two-cell}.

\emph{Step 7 (the case $\Omega = 1$).} Then $\Prb(D = 0) = 0$, and by Step~1
$|\Eop[h\,\ind\{D = 0\}]| \le \Prb(D = 0) = 0$, so that term of Step~4 vanishes and
$\Eop[h] = \Eop[h\,\ind\{D = 1\}] = \Omega\,\Eop[h \mid D = 1] = \Eop[h \mid D = 1]$, the middle
equality being Step~5's definition, available because $\Omega = 1 > 0$. Multiplying by $\Delta_m$
gives $\Dact = M_F$.

\emph{Step 8 (at $\Omega = 1$ the pooled average is not determined by the law).} Step~5's definition
does not apply to $M_P$, since it would divide by $\Prb(D = 0) = 0$. Conditioning on the generated
$\sigma$-algebra supplies no value either. Let $\mathcal D := \sigma(D)$, whose atoms are
$\{D = 1\}$ and $\{D = 0\}$, and let $x$ be any real number. The function
\begin{equation}
\label{eq:version}
  \Delta_m\,\Eop[h \mid D = 1]\,\ind\{D = 1\} \;+\; x\,\ind\{D = 0\}
\end{equation}
is $\mathcal D$-measurable and integrates to $\Delta_m\Eop[h\,\ind_A]$ over every
$A \in \mathcal D$, because two such functions with different values of $x$ differ only on the null
set $\{D = 0\}$. Every real $x$ therefore yields a version of $\Delta_m\Eop[h \mid \mathcal D]$, so
no value of $M_P$ is determined by the law, and the statement at $\Omega = 1$ is Step~7's one-term
equation.

\emph{Step 9 (the case $\Omega = 0$).} Mirror Steps~7 and~8 with the cells exchanged:
$|\Eop[h\,\ind\{D = 1\}]| \le \Prb(D = 1) = 0$, so $\Eop[h] = \Eop[h \mid D = 0]$ and
$\Dact = M_P$, while $M_F$ is left undetermined by the argument of Step~8.
\end{proof}

\medskip
\noindent\emph{The identity is tied to this partition.} The identity \eqref{eq:two-cell} is tied to the disclosure partition and to no other. The flagged
weight factors as $\Omega = \Prb(a = 1)\,\omega_a$ with $\omega_a = \Prb(D = 1 \mid a = 1)$, so
$\Omega \le \Prb(a = 1)$, strictly when $\Prb(a = 1) > 0$ and $\omega_a < 1$. Averaging the kernel
over $\{a = 1\}$ and $\{a = 0\}$ therefore gives a different decomposition, with the engaged but
unflagged histories sitting in the pooled cell of \eqref{eq:two-cell}.
\medskip

\subsubsection*{The factorisation}

Write $\mathcal S = |\partial_\kappa \Dact|$ for the liquidity sensitivity of the premium and
$\mathcal S_P = |\partial_\kappa M_P|$ for the liquidity sensitivity of the pooled cell.

\begin{proposition}[Factorisation of the liquidity sensitivity]
\label{prop:factorisation}
Assume \ref{p:space}--\ref{p:fixed}, the partition of Lemma~\ref{lem:partition}, the two-cell
decomposition of Lemma~\ref{lem:two-cell}, an interior partition $0 < \Omega < 1$ at the policy under
consideration, the invariance of the flagged endpoint in $\kappa$ at fixed policies, and, as an
explicit hypothesis, differentiability of $\kappa \mapsto M_P$. Then
\begin{equation}
\label{eq:factorisation}
  \mathcal S(\kappa,\tau,T) \;=\; \bigl(1 - \Omega(\tau,T)\bigr)\,\mathcal S_P(\kappa,\tau,T),
\end{equation}
exactly. The same factor scales the total variation of $\Dact$ over any grid in $\kappa$, and there
no differentiability is used.
\end{proposition}

\begin{proof}
\emph{Step 1 (the two-cell identity).} The partition is interior at the policy considered, so
Lemma~\ref{lem:two-cell} applies in its non-degenerate branch and \eqref{eq:two-cell} holds at
$(\kappa,\tau,T)$, with $M_F$ and $M_P$ the two cell averages.

% The flagged-cell invariance is the companion result of this subsection ("the flagged cell is
% kappa-free"); the cross-reference is wired when the appendix is assembled.
\emph{Step 2 (the flagged cell does not move with $\kappa$).} At the fixed cutoff and execution
policies of \ref{p:fixed}, the flagged posterior, the flagged price, the entry probability and $M_F$
are invariant to $\kappa$, by the flagged-cell invariance assumed in the statement. Hence
$\partial_\kappa M_F = 0$, and indeed $M_F$ is one and the same number at any two values of
$\kappa$.

\emph{Step 3 (the flagged weight does not move with $\kappa$ either).} This is derived, not assumed.
By Step~4 of Lemma~\ref{lem:partition}, $\{D = 1\}$ is the finite union \eqref{eq:cell-union}, each
of whose ingredients is a function of the signal alone once the cutoff policy is fixed. By
\ref{p:nofeedback} the executed path $B_j(s,d)$ is a function of $(j,s,d)$ alone, so no realised
order flow and no realised price enters it, and with the policies frozen by \ref{p:fixed} the
indicator $D$ is a function of $s$ alone. The law of $s$ carries no $\kappa$: by \ref{p:space} the
signal is $v + \varepsilon$, and by \ref{p:kappa} the laws of $v$ and $\varepsilon$ do not involve
$\kappa$. Indeed $\kappa$ enters the primitives in one place only, the law of the noise mark $z_d$,
which reaches observed pooled order flow $X_d = q_{jd} + z_d$ and reaches nothing that $D$ depends
on. Hence $\Omega = \Prb(D = 1)$ is literally the same number at every $\kappa$, and in particular
$\partial_\kappa\Omega = 0$. This step consumes the freezing of the cutoff vector in \ref{p:fixed},
and no property of the disclosure rule beyond the partition.

\emph{Step 4 (differentiating in $\kappa$).} Two of the three factors on the right of
\eqref{eq:two-cell} are constant in $\kappa$: $M_F$ by Step~2 and $\Omega$ by Step~3. Neither of
those steps uses the present one, so nothing here is circular. With those factors constant,
$\kappa \mapsto \Dact$ is an affine image of $\kappa \mapsto M_P$, which the statement's
differentiability hypothesis makes differentiable; hence $\partial_\kappa\Dact$ exists and equals
$(1 - \Omega)\,\partial_\kappa M_P$. Written with one term per factor that could carry $\kappa$, the
same computation reads
\begin{equation}
\label{eq:three-term}
  \partial_\kappa\Dact
  \;=\; \Omega\,\partial_\kappa M_F
  \;+\; (1 - \Omega)\,\partial_\kappa M_P
  \;+\; (\partial_\kappa\Omega)\,(M_F - M_P),
\end{equation}
the three terms being the flagged cell's own motion, the pooled cell's own motion, and the
reallocation of mass between the cells. Step~2 kills the first and Step~3 the third. The third
deletion assumes nothing about $M_F - M_P$, which is neither small nor signed here; the term goes
because its coefficient is zero.

\emph{Step 5 (the factorisation).} Since $0 < \Omega < 1$ we have $1 - \Omega \in (0,1)$ and
$|1 - \Omega| = 1 - \Omega$. Taking absolute values in the conclusion of Step~4 gives
\eqref{eq:factorisation}, exactly. Step~3 licenses writing the argument of $\Omega$ as $(\tau,T)$
rather than $(\kappa,\tau,T)$, and that convention is kept from here on. Two consequences are worth
recording. First, $\mathcal S > 0$ exactly when $\mathcal S_P > 0$. Second, when $\mathcal S_P > 0$
the map $x \mapsto |x|$ is differentiable at $\partial_\kappa M_P \ne 0$, so $\mathcal S_P$ inherits
whatever differentiability $\partial_\kappa M_P$ has. That transfer creates
no differentiability; it moves it.

\emph{Step 6 (the aggregate form, with no differentiability used).} Fix any grid
$\kappa_0 < \kappa_1 < \dots < \kappa_n$ inside the maintained parameter set, with the partition
interior at each node. Applying Step~1 at $\kappa_i$ and at $\kappa_{i+1}$, and using Steps~2 and~3
in their constancy form, that $M_F$ and $\Omega$ are one and the same number at the two nodes, which
a vanishing derivative at a single $\kappa$ would not deliver, gives
$\Dact(\kappa_{i+1}) - \Dact(\kappa_i) = (1 - \Omega)\bigl[M_P(\kappa_{i+1}) - M_P(\kappa_i)\bigr]$.
Summing absolute values,
\begin{equation}
\label{eq:tv-form}
  \sum_{i=0}^{n-1}\bigl\lvert \Dact(\kappa_{i+1}) - \Dact(\kappa_i)\bigr\rvert
  \;=\; (1 - \Omega)\sum_{i=0}^{n-1}\bigl\lvert M_P(\kappa_{i+1}) - M_P(\kappa_i)\bigr\rvert,
\end{equation}
which is the total-variation clause, finite differences throughout. The same argument runs for any
aggregator of the increment vector that is positively homogeneous of degree one, such as total
variation, mean absolute slope, or the supremum of the absolute increments, because $1 - \Omega$ is
one nonnegative scalar common to every increment. A measurement convention that reports
$\kappa$-sensitivity as a total variation over a grid rather than as a pointwise derivative is then
measuring the same functional.
\end{proof}

% ===========================================================================
% The weight leg at the threshold margin: a tighter threshold raises the
% flagged weight, so the pooled share weakly falls.
% ===========================================================================

\providecommand{\Prb}{\mathbb{P}}
\providecommand{\ind}{\mathbf{1}}

\begin{lemma}[The weight leg at the threshold margin]
\label{lem:threshold-weight}
Fix the liquidity $\kappa$, the window margin $T$, and the plan and cutoff policies. Compare two
threshold margins with $b_0 < \tau' < \tau$, so that $\tau'$ is the tighter threshold, and assume
$\Omega(\tau',T) < 1$. The configuration $b_0 < \tau'$ is what puts the clock equivalence in force
at both thresholds, together with Voice date-monotonicity (standing condition \ref{p:paths}) and
the restriction that only Voice plans cross (standing condition \ref{p:clock}). Then
\begin{enumerate}[label=(\roman*),leftmargin=2.4em,itemsep=3pt]
\item the flagged cells are nested, $\mathcal C_F(\tau,T) \subseteq \mathcal C_F(\tau',T)$, and
every history the tighter threshold newly flags carries the engagement flag $a=1$;
\item the flagged weight rises, $\Omega(\tau',T) \ge \Omega(\tau,T)$, so the weight leg
\begin{equation}
\label{eq:threshold-weight-leg}
W_\tau \;=\; \frac{1-\Omega(\tau',T)}{1-\Omega(\tau,T)}
\end{equation}
is defined and attenuates: $0 < W_\tau \le 1$;
\item if in addition $\Omega(\tau,T) > 0$, $\mathcal S_P(\kappa,\tau,T) > 0$, and the hypotheses of
Proposition~\ref{prop:factorisation} hold at both thresholds, the sensitivity
ratio at this margin factors into the weight leg and the composition leg
$C_\tau = \mathcal S_P(\kappa,\tau',T)/\mathcal S_P(\kappa,\tau,T)$,
\begin{equation}
\label{eq:threshold-ratio}
\frac{\mathcal S(\kappa,\tau',T)}{\mathcal S(\kappa,\tau,T)} \;=\; W_\tau\,C_\tau .
\end{equation}
\end{enumerate}
\end{lemma}

\begin{proof}
Fixed policies hold the map from the signal to the plan, the executed stake path of each plan, and
the law of the signal common to the two rules. The disclosure rule writes the stake path with no
threshold argument, while the crossing date, the filing date, the stake at filing and the
disclosure indicator each carry one. The only objects that move when the threshold moves are
therefore those four, and every comparison below is made at one and the same plan, one and the same
path, and one and the same signal law.

\emph{The product form of the disclosure indicator.} For every plan $j$, every signal $s$, and each
of the two thresholds,
\begin{equation}
\label{eq:threshold-product}
D_j(s;\tau,T) \;=\; \ind\{a_j = 1\}\cdot\ind\{B_j(s,H-T) \ge \tau\} .
\end{equation}
If $a_j = 0$ both sides vanish, the left side because engagement is a conjunct of disclosure and the
right side through its first factor. If $a_j = 1$ the left side is the event that the filing lands
by the horizon: the filing date is the crossing date plus the window margin, and with a finite
window margin the filing landing by the horizon already forces the crossing to be finite, so the
requirement that a crossing occur adds nothing. The clock equivalence then converts the filing
landing by the horizon into the stake $T$ rounds before the horizon having reached the threshold,
which is the second factor. Because $b_0 < \tau' < \tau$, the crossing is the first passage at each
of the two thresholds and the equivalence is available at both.

\emph{The inclusion.} Take any control-node history that the looser threshold flags,
$D_j(s;\tau,T) = 1$. By \eqref{eq:threshold-product} this says $a_j = 1$ and
$B_j(s,H-T) \ge \tau$. Since $\tau' < \tau$, the same real number satisfies
$B_j(s,H-T) \ge \tau > \tau'$, and reading \eqref{eq:threshold-product} at $\tau'$ for the same plan
and the same signal gives $D_j(s;\tau',T) = 1$. The number $B_j(s,H-T)$ that appears in the two
comparisons is one number, and this is where fixed policies are consumed: were the path free to move
with the threshold, the comparison at $\tau'$ would be about a different path. The history was
arbitrary, so $\mathcal C_F(\tau,T) \subseteq \mathcal C_F(\tau',T)$, and taking complements in the
exclusive and exhaustive partition gives
$\mathcal C_P(\tau',T) \subseteq \mathcal C_P(\tau,T)$: a tighter threshold can only shrink the
pooled cell. Every newly flagged history satisfies $D(\tau') = 1$, and by
\eqref{eq:threshold-product} that carries $a = 1$ as a conjunct, so the engagement flag is
identically one on the newly flagged set. Nestedness is a conclusion here rather than a hypothesis.
This is clause~(i).

\emph{The weight leg.} The joint law of the primitives is itself a primitive and does not depend on
the threshold, which enters only through the flagged event. Monotonicity of a probability measure
applied to the inclusion gives
\[
\Omega(\tau',T) \;=\; \Prb\bigl(\mathcal C_F(\tau',T)\bigr)
\;\ge\; \Prb\bigl(\mathcal C_F(\tau,T)\bigr) \;=\; \Omega(\tau,T),
\]
so $1-\Omega(\tau',T) \le 1-\Omega(\tau,T)$. The hypothesis $\Omega(\tau',T) < 1$ makes the
numerator of \eqref{eq:threshold-weight-leg} strictly positive, and the same hypothesis with the
inequality just established gives $\Omega(\tau,T) \le \Omega(\tau',T) < 1$, so the denominator is
strictly positive as well. The quotient is therefore defined, and dividing a weakly smaller strictly
positive number by a strictly positive one gives $0 < W_\tau \le 1$. This is clause~(ii).

\emph{The ratio identity.} With $\Omega(\tau,T) > 0$ the weight-leg inequality puts
both flagged weights strictly inside the unit interval, so Proposition~\ref{prop:factorisation}
applies at each threshold and the noise sensitivity of the expected engagement premium is the
pooled share times the pooled sensitivity,
\[
\mathcal S(\kappa,\tau,T) \;=\; \bigl(1-\Omega(\tau,T)\bigr)\,\mathcal S_P(\kappa,\tau,T),
\qquad
\mathcal S(\kappa,\tau',T) \;=\; \bigl(1-\Omega(\tau',T)\bigr)\,\mathcal S_P(\kappa,\tau',T) .
\]
The flagged weight does not move with $\kappa$, which is what licenses writing it as
$\Omega(\cdot,T)$ with no liquidity argument on either side. Since $1-\Omega(\tau,T) > 0$ and
$\mathcal S_P(\kappa,\tau,T) > 0$, the denominator $\mathcal S(\kappa,\tau,T)$ is strictly positive
and the quotient of the two displays is defined. Dividing and grouping the two factors gives
\eqref{eq:threshold-ratio}, which is clause~(iii).
\end{proof}

\noindent
The weight leg at this margin is signed outright, and it is signed by the geometry of the flagged
set alone: no restriction on the pooled experiment and no restriction on the noise-trading intensity
enters clauses~(i) and~(ii). What clause~(iii) then supplies is the exact form in which the
composition leg has to be signed: $W_\tau C_\tau \le 1$ if and only if $C_\tau \le 1/W_\tau$, and
by clause~(ii) that ceiling is at least one.

% Clock theorem: a shorter clock lowers noise sensitivity iff W_T C_T is at most one.

\begin{theorem}[The clock dial]
\label{thm:clock}
Fix the liquidity $\kappa$, the threshold margin $\tau$, and the plan and cutoff policies. Compare
two window margins $T' < T$, so that $T'$ is the shorter clock. Assume $b_0 < \tau$, that Voice
stake paths are weakly increasing in the calendar date (standing condition \ref{p:paths}), and
that only Voice plans cross (standing condition \ref{p:clock}), which is the configuration that
puts the clock equivalence in force at both windows. Assume $0 < \Omega(\tau,T) < 1$ and
$0 < \Omega(\tau,T') < 1$, and $\mathcal S_P(\kappa,\tau,T) > 0$. Assume that the hypotheses of
Proposition~\ref{prop:factorisation} hold at both clocks. Write
\begin{equation}
\label{eq:clock-legs}
W_T \;=\; \frac{1-\Omega(\tau,T')}{1-\Omega(\tau,T)},
\qquad
C_T \;=\; \frac{\mathcal S_P(\kappa,\tau,T')}{\mathcal S_P(\kappa,\tau,T)} .
\end{equation}
Then
\begin{enumerate}[label=(\roman*),leftmargin=2.4em,itemsep=3pt]
\item the weight leg attenuates: $0 < W_T \le 1$;
\item the sensitivity ratio factors into the two legs,
\begin{equation}
\label{eq:clock-ratio}
\frac{\mathcal S(\kappa,\tau,T')}{\mathcal S(\kappa,\tau,T)} \;=\; W_T\,C_T ;
\end{equation}
\item the shorter clock lowers noise sensitivity exactly when the product of the two legs is at
most one,
\begin{equation}
\label{eq:clock-iff}
\mathcal S(\kappa,\tau,T') \;\le\; \mathcal S(\kappa,\tau,T)
\qquad\Longleftrightarrow\qquad
W_T\,C_T \;\le\; 1 .
\end{equation}
\end{enumerate}
\end{theorem}

\begin{proof}
Fixed policies hold the map from the signal to the plan, the executed stake path of each plan, and
the law of the signal common to the two rules, so every object below is compared at one and the
same plan and one and the same signal law.

\emph{The weight leg.} Let $j$ be a Voice plan and $s$ a signal whose history is flagged under the
longer clock, $D_j(s;\tau,T) = 1$. By the clock equivalence, the filing lands by the horizon exactly
when the stake $T$ rounds before the horizon has reached the threshold, so
$B_j(s,H-T) \ge \tau$. Since $T' < T$ gives $H-T' > H-T$, and since the stake path of a Voice plan
is weakly increasing in the round, $B_j(s,H-T') \ge B_j(s,H-T) \ge \tau$. The engagement flag
$a_j = 1$ does not move with the clock. Reading the clock equivalence in the other direction at the
shorter clock gives $D_j(s;\tau,T') = 1$. Only Voice plans reach the threshold, so no other history
has to be checked, and the flagged cell at the longer clock is contained in the flagged cell at the
shorter clock. Taking probabilities under the common signal law, $\Omega(\tau,T') \ge
\Omega(\tau,T)$, so $1 - \Omega(\tau,T') \le 1 - \Omega(\tau,T)$. Both sides are strictly
positive because $\Omega < 1$ at each clock, so the quotient in \eqref{eq:clock-legs} is defined and
$0 < W_T \le 1$. This is clause~(i).

\emph{The ratio identity.} By Proposition~\ref{prop:factorisation} at both clocks, which the
statement assumes, the noise sensitivity of the expected engagement premium is the pooled share
times the pooled sensitivity. The factorisation hypotheses at the shorter clock, differentiability
of $\kappa \mapsto M_P(\kappa,\tau,T')$ included, make $\mathcal S_P(\kappa,\tau,T')$ exist, so
$C_T$ in \eqref{eq:clock-legs} is defined. At each clock,
\[
\mathcal S(\kappa,\tau,T) \;=\; \bigl(1-\Omega(\tau,T)\bigr)\,\mathcal S_P(\kappa,\tau,T),
\qquad
\mathcal S(\kappa,\tau,T') \;=\; \bigl(1-\Omega(\tau,T')\bigr)\,\mathcal S_P(\kappa,\tau,T') .
\]
The flagged weight does not move with $\kappa$, which is what licenses writing it as
$\Omega(\tau,\cdot)$ with no liquidity argument on either side. Since $1-\Omega(\tau,T) > 0$ and
$\mathcal S_P(\kappa,\tau,T) > 0$, the denominator $\mathcal S(\kappa,\tau,T)$ is strictly positive
and the quotient of the two displays is defined. Dividing and grouping the two factors gives
\eqref{eq:clock-ratio}, which is clause~(ii).

\emph{The equivalence.} Multiplying an inequality by the strictly positive number
$\mathcal S(\kappa,\tau,T)$ preserves it in both directions, so
$\mathcal S(\kappa,\tau,T') \le \mathcal S(\kappa,\tau,T)$ holds exactly when
$\mathcal S(\kappa,\tau,T')/\mathcal S(\kappa,\tau,T) \le 1$, and by \eqref{eq:clock-ratio} that
quotient is $W_T C_T$. This is \eqref{eq:clock-iff} and clause~(iii).
\end{proof}

\noindent
Because $W_T > 0$, the criterion \eqref{eq:clock-iff} can be read as a ceiling on the composition
leg: the shorter clock lowers noise sensitivity exactly when
\[
C_T \;\le\; \frac{1}{W_T} \;=\; \frac{1-\Omega(\tau,T)}{1-\Omega(\tau,T')} .
\]
By clause~(i) that ceiling is at least one, so $C_T \le 1$ is sufficient for attenuation at the
clock margin.


% proofs/08_blackwell.tex
% The tagged experiment under a tighter disclosure rule.
% Fragment, to be \input by appendix.tex.

\section{Information is monotone in the rule}
\label{sec:blackwell}

Let $\theta=(j(s),s)$ and fix a rule $r=(\tau,T)$. At date $d$, set
$D_r^d=\ind\{a=1,\ c(\theta;\tau)+T\leq d\}$ and $L_{r,e}=D_r^e$. At the control date the market
observes the tagged output
\[
Y_r=
\begin{cases}
(F,\mathsf S_{F,r}),&D_r^H=1,\\
(P,\mathcal H^P_r),&D_r^H=0,
\end{cases}
\qquad
\mathsf S_{F,r}=(B_r^F,Q_r^F,a{=}1),
\]
where
\[
\mathcal H^P_r=(X_0,\ldots,X_H;L_{r,0},\ldots,L_{r,H}).
\]
Write $\mathcal C_{F,r}=\{D_r^H=1\}$, $\mathcal C_{P,r}=\{D_r^H=0\}$, and
$\Omega_r=\Prb(D_r^H=1)$, used below in the strictness criterion.
Every flag coordinate in a pooled history is zero. Prices may be appended to either branch. They
are pinned measurable functions of the displayed information and do not enlarge the experiment.
Write $\mathcal E_r(\kappa)$ for the conditional channel from $\theta$ to $Y_r$ at liquidity
$\kappa$.

\csname unlabelledtrue\endcsname
\begin{theorem}[Tightening either disclosure dial]
\label{thm:blackwell}
Fix a common $\kappa\in[0,1]$ and fixed policies. Assume:
\begin{enumerate}[label=(B\arabic*),leftmargin=3.0em,itemsep=3pt]
\item \emph{Standing conditions.} Conditions \ref{p:space}--\ref{p:decoder} hold. Condition (S14),
  identified flagged types, is load-bearing: on the tighter rule's
  flagged region the tuple recovers the signal, and hence the type.
\item \emph{A tighter rule.} Either $b_0<\tau'<\tau$ at a common $T$, with
  $r_+=(\tau',T)$ and $r_-=(\tau,T)$, or $T'<T$ at a common $\tau$, with
  $r_+=(\tau,T')$ and $r_-=(\tau,T)$.
\end{enumerate}
Then the tighter tagged experiment is weakly more informative in Blackwell's order about
$\theta$ and about the price-relevant state $(v,a)$. Specifically, there is a Markov kernel $K_H$,
independent of the true type, such that
\begin{equation}
\label{eq:blackwell-kernel}
\mathcal L(Y_{r_-}\mid\theta)
=\int K_H(y,\mathord\cdot)\,\mathcal L(Y_{r_+}\in\mathrm dy\mid\theta).
\end{equation}
Thus $\mathcal E_{r_+}(\kappa)\succeq_B\mathcal E_{r_-}(\kappa)$ pointwise under the pointwise
form of (S14), or modulo a prior-null set under its almost-sure form. No positive-mass
condition on either cell is needed. The endpoint values of $\kappa$ refer to the primitive noise
law in (S12).
\end{theorem}
\csname unlabelledfalse\endcsname

\begin{proof}
First, the cells are nested. At the threshold margin, a path that crosses $\tau$ crosses
$\tau'<\tau$ no later. At the clock margin, $c+T\leq H$ implies $c+T'\leq H$. Conditions (S4),
(S5), (S7), and (S11) therefore give
\begin{equation}
\label{eq:blackwell-nesting}
\mathcal C_{F,r_-}\subseteq\mathcal C_{F,r_+},
\qquad
\mathcal C_{P,r_+}\subseteq\mathcal C_{P,r_-}.
\end{equation}

On $\mathcal C_{F,r_+}$, condition (S14) gives
$B^F_{r_+}+Q^F_{r_+}=g(s)$ with $g$ strictly increasing. Its inverse on its image is measurable,
and the Lusin--Suslin theorem makes the image Borel, so $\mathsf S_{F,r_+}$ recovers $s$, then
$j(s)$, and hence $\theta$. Denote this decoder by
$\iota_+$.

Define $K_H$ on the three branches of the tighter output. If $y=(P,\mathfrak h)$, emit
$(P,\mathfrak h)$. If $y=(F,\mathsf s_F)$, recover
$\widehat\theta=\iota_+(\mathsf s_F)$. When $D_{r_-}^H(\widehat\theta)=1$, emit the deterministic
looser tuple $(F,\mathsf S_{F,r_-}(\widehat\theta))$. When
$D_{r_-}^H(\widehat\theta)=0$, draw $z'_{0:H}\sim Q_\kappa$ and emit
\[
\left(P,
\bigl(m_{r_-,0}(\widehat\theta)+z'_0,\ldots,m_{r_-,H}(\widehat\theta)+z'_H;
0,\ldots,0\bigr)\right).
\]
Here $m_{r_-,d}$ denotes the looser rule's mark path.
Define the kernel arbitrarily on flagged tuples outside the image of the tighter flagged channel.
The decoder, disclosure test, tuples, and marks are measurable by (S1) to (S7) and (S14), and the
last branch uses the fixed law in (S12), so this is a Markov kernel that reads only $y$.

Fix a type. If it is pooled under the tighter rule, nesting makes it pooled under the looser rule.
The first branch is exact because (S7), (S11), and (S12) give the same mark path, noise draw, and
all-zero flag coordinates under both rules. If the type is flagged under both, the second branch
emits its looser tuple. If it is newly flagged, the third branch redraws the common noise channel
and adds it to the looser mark path, which is exactly its looser pooled law. These cases prove
\eqref{eq:blackwell-kernel}. If prices are recorded, the kernel discards the input prices and
appends the pinned prices belonging to its emitted output.

Finally, $a$ is a function of $\theta$, the conditional law of $v$ given $\theta$ is
rule-invariant, and the output is conditionally independent of $v$ given $\theta$ by (S7) and
(S12). Integrating \eqref{eq:blackwell-kernel} against the common conditional law of $\theta$
given $(v,a)$ proves the same order about $(v,a)$.
\end{proof}

\csname unlabelledtrue\endcsname
\begin{corollary}[Posterior engagement risk]
\label{cor:blackwell-posterior}
Assume:
\begin{enumerate}[label=(P\arabic*),leftmargin=3.0em,itemsep=3pt]
\item \emph{Standing conditions.} Conditions (S1) to (S14) hold.
\item \emph{Rule comparison.} The two rules satisfy (B2) at a common $\kappa$, so
  Theorem~\ref{thm:blackwell} applies.
\end{enumerate}
Then the expected posterior risk of engagement, measured by posterior variance, falls weakly under
the tighter rule:
\begin{equation}
\label{eq:blackwell-posterior}
\Eop\!\left[\operatorname{Var}(a\mid Y_{r_+})\right]
\leq
\Eop\!\left[\operatorname{Var}(a\mid Y_{r_-})\right].
\end{equation}
The same inequality holds with $a$ replaced by any square-integrable function of $\theta$.
\end{corollary}
\csname unlabelledfalse\endcsname

\begin{proof}
Use $K_H$ to couple the outputs so that $\theta$, $Y_{r_+}$, and $Y_{r_-}$ form a Markov chain in
that order. Put $M_+=\Eop[a\mid Y_{r_+}]$ and $M_-=\Eop[a\mid Y_{r_-}]$. Then
$M_-=\Eop[M_+\mid Y_{r_-}]$. Conditional variance decomposition gives
\[
\Eop\!\left[\operatorname{Var}(a\mid Y_{r_-})\right]
-\Eop\!\left[\operatorname{Var}(a\mid Y_{r_+})\right]
=\Eop\!\left[\operatorname{Var}(M_+\mid Y_{r_-})\right]\geq0.
\]
The same argument applies to every square-integrable function of the type.
\end{proof}

\csname unlabelledtrue\endcsname
\begin{corollary}[The two-parameter Blackwell chain]
\label{cor:blackwell-chain}
Assume:
\begin{enumerate}[label=(K\arabic*),leftmargin=3.0em,itemsep=3pt]
\item \emph{Standing conditions.} Conditions (S1) to (S14) hold at both rules.
\item \emph{Rule comparison.} The pair $r_+,r_-$ satisfies (B2).
\item \emph{Liquidity comparison.} The active order is two noise lumps and
  $0<\kappa_L<\kappa_H<1$, so Lemma~\ref{lem:g2} applies on each pooled cell.
\end{enumerate}
Then
\begin{equation}
\label{eq:blackwell-chain}
\mathcal E_{r_+}(\kappa_L)
\succeq_B\mathcal E_{r_+}(\kappa_H)
\succeq_B\mathcal E_{r_-}(\kappa_H),
\end{equation}
and also
\[
\mathcal E_{r_+}(\kappa_L)
\succeq_B\mathcal E_{r_-}(\kappa_L)
\succeq_B\mathcal E_{r_-}(\kappa_H).
\]
In particular, a tighter rule at lower liquidity dominates a looser rule at higher liquidity.
\end{corollary}
\csname unlabelledfalse\endcsname

\begin{proof}
Fix one rule. Its cell event and flagged tuple do not move with $\kappa$ under (S7), (S10),
(S11), and (S12). Apply the kernel of Lemma~\ref{lem:g2} on the pooled branch and leave the flagged
tuple unchanged. This garbles the tagged experiment at $\kappa_L$ into the one at
$\kappa_H$. If prices are appended, the pooled price is recomputed from the emitted history and
the flagged price is recomputed from the emitted tuple. Theorem~\ref{thm:blackwell}
gives the rule order at either fixed liquidity. Transitivity gives both displayed chains.
\end{proof}

\csname unlabelledtrue\endcsname
\begin{corollary}[When the Blackwell order is strict]
\label{cor:blackwell-strict}
Assume:
\begin{enumerate}[label=(R\arabic*),leftmargin=3.0em,itemsep=3pt]
\item \emph{Standing conditions.} Conditions (S1) to (S14) and the comparison in (B2) hold at a
  common $\kappa\in[0,1]$.
\item \emph{Positive new flagged mass.} The set
  $N=\mathcal C_{F,r_+}\setminus\mathcal C_{F,r_-}$ has positive mass. By the nesting in
  Theorem~\ref{thm:blackwell}, $\Prb(N)=\Omega_{r_+}-\Omega_{r_-}$.
\item \emph{A nonatomic signal on the cut.} The conditional law of $s$ given $N$ is nonatomic.
\item \emph{A finite pooled range.} The looser pooled history takes values in a finite set.
\end{enumerate}
Then the order about $\theta$ in Theorem~\ref{thm:blackwell} is strict: the looser experiment cannot
garble the tighter one. \textsc{Numerical.} On the calibration's threshold ladder, these conditions
hold for each positive-mass cut at $T=5$, so every such threshold comparison is strict. The
threshold cuts at $T=10$ are null, and the corresponding rule experiments are
Blackwell-equivalent. These calibration facts come from
\texttt{numerical\_v4/checks/t2\_threshold\_revelation\_check.json}.
\end{corollary}
\csname unlabelledfalse\endcsname

\begin{proof}
Theorem~\ref{thm:blackwell} gives the weak order. Suppose a reverse kernel existed. By (S7) and
(S13), for every type in $N$, the tighter output is the point mass at its flagged tuple, while the
looser output lies in
the finite pooled range. A mixture of rows of the reverse kernel can equal that point mass only if
every row receiving positive weight is the same point mass. The tighter flagged tuple would
therefore take only finitely many values on $N$ outside the type-null exceptional set on which the
reverse-kernel identity need not hold. Deleting that set preserves $\Prb(N)>0$ and the nonatomic
conditional signal law. Condition (S14) makes that tuple
injective in $s$, while (R3) makes its conditional distribution on $N$ nonatomic. This is a
contradiction, so no reverse kernel exists. The final calibration sentences apply this criterion to
the positive cuts at $T=5$. At the reported null cuts, $T=10=H$ makes every flagged type cross at
date zero, and nesting plus nullity makes the flagged sets coincide up to a prior-null set. On that
set $B^F=b^*(s)$ and $Q^F=0$ under both thresholds, so identity kernels on the flagged and pooled
branches give the stated equivalence.
\end{proof}

\begin{remark}[What the order does not sign]
A Blackwell order at fixed liquidity does not sign the liquidity sensitivity of the engagement
premium $\Delta_m\Eop[h(\mathcal I_H)]$, where $h=\pi p$. It compares level experiments at one
$\kappa$, while changing the rule also changes the type composition and weight of the pooled cell
whose derivative defines that sensitivity.
\end{remark}


% =============================================================================
% proofs/02_garbling.tex
% The garbling lemma at order size two and the threshold dial.
% Statements first, then the proofs.  Assembled into appendix.tex.
% Needs amsmath, amssymb, amsthm and enumitem in the assembling preamble.
% =============================================================================

\providecommand{\Eop}{\mathbb{E}}
\providecommand{\Prb}{\mathbb{P}}
\providecommand{\ind}{\mathbf{1}}
\providecommand{\Dact}{\Delta^{\mathrm{act}}}
\providecommand{\Wrev}{\mathcal{W}}
\providecommand{\Grev}{\mathcal{G}}

\makeatletter
\@ifundefined{c@theorem}{\newtheorem{theorem}{Theorem}}{}
\@ifundefined{c@lemma}{\newtheorem{lemma}{Lemma}}{}
\@ifundefined{c@condition}{\newtheorem{condition}{Condition}}{}
\@ifundefined{c@remark}{\newtheorem{remark}{Remark}}{}
\makeatother

\section{Order size two: the pooled experiment and the threshold dial}
\label{sec:pf-garbling}

\subsection*{Setting and notation}
\label{sec:pf-garbling-setup}

Trading runs over rounds $d = 0,\ldots,H$. Noise marks $(z_d)_{d=0}^H$ are independent and
identically distributed across rounds and independent of the blockholder's type $\theta$ (and thus
of the signal $s$ and mark path $(m_d)_{d=0}^H$). In each round, noise takes the values $-\bar z$,
$0$ and $+\bar z$ with probabilities $\kappa/2$, $1-\kappa$ and $\kappa/2$. The blockholder's order
while she is building the stake is two noise lumps, so the pooled order mark is
$q_{jd}(s) \in \{0, 2\bar z\}$ and pooled order flow is $X_d = q_{jd}(s) + z_d$. The support of $X_d$ is
$\{-\bar z, 0, \bar z, 2\bar z, 3\bar z\}$: a round in which the blockholder buys, an
\emph{active} round, produces $X_d \in \{\bar z, 2\bar z, 3\bar z\}$, and a round in which she
does not, an \emph{idle} round, produces $X_d \in \{-\bar z, 0, \bar z\}$. The two supports meet
at $\bar z$ and nowhere else. Write
\begin{equation}
\label{eq:g-eps}
\varepsilon \;:=\; \kappa/2 \;\in\;(0,\tfrac12) \qquad\text{for }\kappa\in(0,1).
\end{equation}

Policies are fixed throughout: the plan menu, the cutoff vector and the execution paths are held
at one setting while the disclosure rule $(\tau,T)$ moves. Let $\theta$ denote the blockholder's
type, the pair of plan and signal, and let
\[
m_d(\theta) \;\in\; \{0,\,2\bar z\}, \qquad d = 0,\ldots,H,
\]
be its mark path, equal to $2\bar z$ on the rounds in which the plan buys and $0$ elsewhere. The
mark path is a deterministic function of the type and takes finitely many values. The disclosure
indicator $D(\theta) = \ind\{a(\theta) = 1,\ c(\theta;\tau) + T \le H\}$ is a function of the type
alone, because the stake path does not respond to realised order flow and the crossing date reads
off that path. Write $\rho$ for the prior law of $\theta$, $\Omega(\tau,T) = \Prb(D=1)$ for the
flagged weight, and, when $\Omega(\tau,T) < 1$,
\[
\rho_P^{\tau,T} \;:=\; \rho(\,\cdot \mid D = 0\,)
\]
for the pooled type law, which does not depend on $\kappa$.

The kernel is $h(\mathcal I) = \pi(\mathcal I)\,p(\mathcal I)$, with $\pi$ the engagement
posterior and $p$ the entry probability. The price and the entry probability at an information
set are pinned by the posterior engagement share $\pi$ and the posterior mean $\hat v$ of the
fundamental, so
\begin{equation}
\label{eq:g-kernel}
h(\nu) \;=\; \mathsf h\bigl(\pi(\nu),\,\hat v(\nu)\bigr),
\end{equation}
where $\nu$ is the posterior over types and $\pi(\nu)$, $\hat v(\nu)$ are both linear in $\nu$.
The map $\mathsf h$ itself carries no $\kappa$: $\kappa$ enters $h(\nu)$ only through the
posterior $\nu$. Lemma~\ref{lem:g3}(a) rests on that, treating each $\Grev(S)$ as a
$\kappa$-free number. A curvature statement about $h$ is read as a statement about $\mathsf h$
on the convex hull of the pairs $(\pi,\hat v)$ the pooled cell generates. The premium objects are
$\Dact = \Delta_m \Eop[h(\mathcal I_H)]$, $M_F = \Delta_m \Eop[h \mid D=1]$,
$M_P = \Delta_m \Eop[h \mid D=0]$, and the sensitivities are
$\mathcal S = |\partial_\kappa \Dact|$ and $\mathcal S_P = |\partial_\kappa M_P|$. For a
threshold pair $\tau' < \tau$ at a common window,
\[
W_\tau \;=\; \frac{1-\Omega(\tau',T)}{1-\Omega(\tau,T)},
\qquad
C_\tau \;=\; \frac{\mathcal S_P(\tau',T)}{\mathcal S_P(\tau,T)} .
\]

For a set of rounds $S$, the \emph{$S$-record} is the pair
$\bigl(S,\ (m_e(\theta))_{e \in S}\bigr)$; Lemma~\ref{lem:g1}(b) shows that the pooled history
determines it. Finally, call an event $\mathcal A$ a \emph{cell event} when it is determined by
the type alone, so that conditioning on it leaves the noise law untouched. The pooled cell
$\{D = 0\}$ is a cell event, and so is the event that the filing has not landed by a date $d$.
For a cell event $\mathcal A$ with $\Prb(\mathcal A) > 0$, write $\rho_{\mathcal A} := \rho(\,\cdot \mid \mathcal A\,)$.

\subsection*{Statements}
\label{sec:pf-garbling-statements}

\begin{lemma}[Erasure form of the pooled experiment]
\label{lem:g1}
Fix a policy $(\tau,T)$, a depth $d \le H$, and $\kappa \in (0,1)$. Assume the noise marks
$(z_e)_{e \le d}$ are i.i.d.\ and independent of $\theta$. Let
\[
R_d \;:=\; \{\, e \le d \;:\; X_e \neq \bar z \,\}
\]
be the set of rounds whose flow differs from one lump. Then:
\begin{enumerate}[label=(\alph*),leftmargin=2.4em,itemsep=3pt]
\item $\Prb(X_e = \bar z \mid \theta) = \varepsilon$ for every type and every round, so the value
  $\bar z$ carries likelihood ratio one across types and its probability rises with $\kappa$;
  consequently $R_d$ is independent of $\theta$ and each round belongs to it independently with
  probability $1-\varepsilon$;
\item on $\{R_d = S\}$ the flow determines $m_e(\theta)$ for every $e \in S$, and conditional on
  the $S$-record the law of $X_{0:d}$ does not depend on $\theta$;
\item consequently the $S$-record is sufficient: for every cell event $\mathcal A$ with
  $\Prb(\mathcal A) > 0$, the posterior after observing $X_{0:d}$ on $\mathcal A$ is
  \[
  \nu_{0:d} \;=\; \rho_{\mathcal A}\bigl(\,\cdot \mid (m_e)_{e \in R_d}\,\bigr),
  \]
  and in particular, when $\Omega(\tau,T) < 1$, the pooled posterior at the control node takes
  this form with $\mathcal A = \{D=0\}$ and $d = H$.
\end{enumerate}
\end{lemma}

\begin{lemma}[Liquidity garbles the pooled experiment]
\label{lem:g2}
Fix a policy and a depth $d \le H$, let $0 < \kappa < \kappa' < 1$, and assume the noise marks
$(z_e)_{e \le d}$ are i.i.d.\ and independent of $\theta$. There is a Markov kernel
$\Lambda$ from histories to histories, not depending on the type, such that for every type
$\theta$ the image under $\Lambda$ of the law of $X_{0:d}$ at $\kappa$ is its law at $\kappa'$.
One such kernel deletes each element of $R_d$ independently with probability
$(\kappa'-\kappa)/(2-\kappa)$, emits $\bar z$ on every deleted round and on every round outside
$R_d$, and redraws the flow on the surviving rounds from the $\kappa'$ conditional law given the
mark. The pooled experiment at $\kappa'$ is therefore a garbling of the pooled experiment at
$\kappa$, and the law of the pooled posterior at $\kappa'$ is a mean preserving contraction of
its law at $\kappa$, and when $\Omega(\tau,T) < 1$, both have mean $\rho_P^{\tau,T}$.
\end{lemma}

\begin{lemma}[Exact liquidity representation of the pooled premium]
\label{lem:g3}
Fix a policy $(\tau,T)$, a cell event $\mathcal A$ of positive prior mass, a depth $d \le H$, and
assume the noise marks $(z_e)_{e \le d}$ are i.i.d.\ and independent of $\theta$.
For $S \subseteq \{0,\ldots,d\}$ define
\begin{equation}
\label{eq:g-G}
\Grev^{\,\mathcal A, d}(S) \;:=\;
\Eop_{\rho_{\mathcal A}}\Bigl[\, h\bigl(\rho_{\mathcal A}(\,\cdot \mid (m_e)_{e\in S})\bigr)
\,\Bigr],
\end{equation}
a number that does not depend on $\kappa$. The three parts below are written for the pooled cell
at the control node, $\mathcal A = \{D=0\}$ and $d = H$, where
$\Grev_{\tau,T} := \Grev^{\,\{D=0\},H}$ and $\Omega(\tau,T) < 1$; each holds verbatim for any
cell event and any depth $d \le H$, with $H$ replaced by $d$ throughout.
\begin{enumerate}[label=(\alph*),leftmargin=2.4em,itemsep=3pt]
\item \emph{Representation.} For every $\kappa \in (0,1)$,
  \begin{equation}
  \label{eq:g-rep}
  M_P(\tau,T;\kappa) \;=\; \Delta_m \sum_{S \subseteq \{0,\ldots,H\}}
  (1-\varepsilon)^{|S|}\,\varepsilon^{\,H+1-|S|}\; \Grev_{\tau,T}(S).
  \end{equation}
\item \emph{Derivative.} With the $\kappa$-free coefficients
  \begin{equation}
  \label{eq:g-ck}
  c_k(\tau,T) \;:=\; \sum_{d=0}^{H}\ \
  \sum_{\substack{S \subseteq \{0,\ldots,H\}\setminus\{d\} \\ |S| = k}}
  \Bigl[\,\Grev_{\tau,T}(S \cup \{d\}) - \Grev_{\tau,T}(S)\,\Bigr],
  \qquad k = 0,\ldots,H,
  \end{equation}
  the pooled premium is differentiable in $\kappa$ on $(0,1)$ and
  \begin{equation}
  \label{eq:g-deriv}
  \partial_\kappa M_P(\tau,T;\kappa) \;=\; -\,\frac{\Delta_m}{2}\,\Wrev_{\tau,T}(\kappa),
  \qquad
  \Wrev_{\tau,T}(\kappa) \;:=\; \sum_{k=0}^{H}(1-\varepsilon)^k \varepsilon^{\,H-k}\,c_k(\tau,T).
  \end{equation}
\item \emph{Sign.} If $\mathsf h$ of \eqref{eq:g-kernel} is concave on the convex hull of the
  pairs $(\pi,\hat v)$ the pooled cell generates, then every increment in \eqref{eq:g-ck} is at
  most zero, hence $c_k(\tau,T)\le 0$ for every $k$ and $\partial_\kappa M_P \ge 0$ at every
  $\kappa \in (0,1)$. If $\mathsf h$ is convex there, all three inequalities reverse. The pooled
  expectation of the kernel is therefore monotone in $\kappa$, in the direction the sign of the
  chord of $\mathsf h$ gives, and the same holds at every depth $d \le H$.
\end{enumerate}
\end{lemma}

\begin{condition}[Revelation dominance at a threshold pair]
\label{cond:D}
Let $b_0 < \tau' < \tau$ at a common window $T$, and let $\Wrev$ be as in \eqref{eq:g-deriv}. The
pair satisfies revelation dominance at $\kappa$ if
\begin{equation}
\label{eq:g-D}
\bigl|\Wrev_{\tau',T}(\kappa)\bigr| \;\le\; \bigl|\Wrev_{\tau,T}(\kappa)\bigr| .
\end{equation}
\end{condition}

\csname unlabelledtrue\endcsname
\par\noindent\textsc{Labels: PROVED and NUMERICAL.} Parts~(A) and~(B) and the implication in
part~(C) are PROVED; Condition~\ref{cond:D}, and with it the conclusion of part~(C), is NUMERICAL
on the grid Remark~\ref{rem:g-grid} names.\par\nopagebreak
\begin{theorem}[The threshold dial at fixed policies]
\label{thm:g-threshold}
Fix the plan and cutoff policies and a window $T$, and take $b_0 < \tau' < \tau$ with
$\Omega(\tau,T) > 0$, $\Omega(\tau',T) < 1$ and $\mathcal S_P(\tau,T) > 0$; part~(B) then puts
$\Omega(\tau,T) \le \Omega(\tau',T) < 1$, so both pooled cells carry positive mass. Assume the
two cell decomposition $\Dact = \Omega M_F + (1-\Omega)M_P$ and the $\kappa$-invariance of the
flagged average $M_F$, each with its own hypotheses. Then:
\begin{enumerate}[label=(\Alph*),leftmargin=2.4em,itemsep=3pt]
\item \emph{Factorisation.} $\mathcal S = (1-\Omega)\,\mathcal S_P$, exactly, at every
  $\kappa \in (0,1)$.
\item \emph{Weight leg.} $\mathcal C_F(\tau,T) \subseteq \mathcal C_F(\tau',T)$ and every newly
  flagged history is generated by a Voice plan, so $\Omega(\tau',T)\ge\Omega(\tau,T)$ and
  $W_\tau \in [0,1]$, with equality when no history is reclassified (that is, up to a null set
  of histories). This leg carries no condition beyond the standing ones.
\item \emph{Composition leg and the dial.} Under Condition~\ref{cond:D} at the pair and at
  $\kappa$, $C_\tau \le 1$, and therefore
  \begin{equation}
  \label{eq:g-dial}
  \frac{\mathcal S(\tau',T)}{\mathcal S(\tau,T)} \;=\; W_\tau\,C_\tau \;\le\; 1 .
  \end{equation}
  Equality holds when no history is reclassified (up to a null set of histories), since the two
  rules then induce the same partition almost surely and both factors are one.
\end{enumerate}
\end{theorem}
\csname unlabelledfalse\endcsname

\begin{remark}[A $\kappa$-free sufficient form of Condition~\ref{cond:D}]
\label{rem:g-Dstar}
The coefficients $c_k(\tau,T)$ of \eqref{eq:g-ck} do not depend on $\kappa$, so
Condition~\ref{cond:D} is a comparison of two vectors in $\mathbb R^{H+1}$. If the $H+1$ numbers
$c_0(\tau,T),\ldots,c_H(\tau,T)$ share one sign and
\begin{equation}
\label{eq:g-Dstar}
|c_k(\tau',T)| \;\le\; |c_k(\tau,T)| \qquad\text{for every } k = 0,\ldots,H,
\end{equation}
then Condition~\ref{cond:D} holds at every $\kappa \in (0,1)$ at once.
\end{remark}

\begin{remark}[What the coefficients are]
\label{rem:g-cells}
By Lemma~\ref{lem:g1} the $S$-record partitions the pooled type set into the level sets of the
restricted mark path $\theta \mapsto (m_e(\theta))_{e \in S}$, so $\Grev_{\tau,T}(S)$ is a finite
sum over those level sets of the cell mass times the kernel at the cell posterior. Each
$c_k(\tau,T)$ is therefore a closed form expression in the pooled type weights, computable
without reference to $\kappa$ and without enumerating order flow.
\end{remark}

\subsection*{Proofs}
\label{sec:pf-garbling-proofs}

\begin{proof}[Proof of Lemma~\ref{lem:g1}]
Condition on a type $\theta$ and a round $e \le d$. If the round is active the mark is $2\bar z$
and $X_e$ takes the values $\bar z$, $2\bar z$, $3\bar z$ with probabilities $\kappa/2$,
$1-\kappa$, $\kappa/2$. If the round is idle the mark is zero and $X_e$ takes the values
$-\bar z$, $0$, $\bar z$ with the same three probabilities in the same order. In both cases
\[
\Prb(X_e = \bar z \mid \theta) \;=\; \kappa/2 \;=\; \varepsilon ,
\]
which proves the first claim of part~(a); the noise draws are independent across rounds and
independent of the type, so the indicators $\ind\{e \in R_d\}$, $e \le d$, are independent
Bernoulli variables with success probability $1-\varepsilon$, independent of $\theta$. That is
part~(a).

For part~(b), the two supports meet only at $\bar z$: an active round produces
$X_e \in \{2\bar z, 3\bar z\}$ and an idle round produces $X_e \in \{-\bar z, 0\}$ whenever
$X_e \neq \bar z$. So on $\{e \in R_d\}$ the flow determines the mark. It remains to check that
nothing else is carried. Conditional on the round being active and revealing,
\[
\Prb(X_e = 2\bar z \mid \text{active},\, e \in R_d) = \frac{1-\kappa}{1-\varepsilon},
\qquad
\Prb(X_e = 3\bar z \mid \text{active},\, e \in R_d) = \frac{\varepsilon}{1-\varepsilon},
\]
and conditional on the round being idle and revealing the values $0$ and $-\bar z$ carry those
same two probabilities. The two conditional laws are the same pair of numbers, so once the mark
is given, which of its two revealing values was realised is a draw whose law does not depend on
the type. Rounds outside $R_d$ carry the single value $\bar z$ and no residual randomness at all.
Combining across rounds by independence, the conditional law of $X_{0:d}$ given the $S$-record
does not depend on $\theta$, which is part~(b).

Part~(c) is the sufficiency conclusion. Write $\mathcal R$ for the $S$-record. Since the
conditional law of $X_{0:d}$ given $(\mathcal R,\theta)$ does not depend on $\theta$,
$\theta$ and $X_{0:d}$ are conditionally independent given $\mathcal R$, and Bayes' rule gives
$\Prb(\theta \in \cdot \mid X_{0:d}, D=0) = \Prb(\theta \in \cdot \mid \mathcal R, D=0)$. The
event $\{D=0\}$ is a type event, so conditioning on it is conditioning $\rho$ down to
$\rho_P^{\tau,T}$, and the displayed posterior is the one stated.
\end{proof}

\begin{proof}[Proof of Lemma~\ref{lem:g2}]
Let $\varepsilon = \kappa/2 < \varepsilon' = \kappa'/2$ and put
\[
\delta \;:=\; \frac{\varepsilon' - \varepsilon}{1-\varepsilon}
\;=\; \frac{\kappa' - \kappa}{2 - \kappa} \;\in\;(0,1).
\]
Define $\Lambda$ on a $\kappa$-history as follows. First read off the record
$(R_d, (m_e)_{e \in R_d})$, which Lemma~\ref{lem:g1}(b) allows without knowing the type. Second,
delete each element of $R_d$ independently with probability $\delta$, obtaining $R'_d$. Third,
emit the flow value $\bar z$ on every round outside $R'_d$, and on every round $e \in R'_d$ emit
one of the two revealing values consistent with $m_e$, drawn with probabilities
$(1-\kappa')/(1-\varepsilon')$ for the zero noise value and $\varepsilon'/(1-\varepsilon')$ for
the other. The kernel reads only the history, never the type.

Fix a type. Under the $\kappa$ law each round lies in $R_d$ independently with probability
$1-\varepsilon$, so under $\Lambda$ each round lies in $R'_d$ independently with probability
\[
(1-\varepsilon)(1-\delta) \;=\; (1-\varepsilon) - (\varepsilon'-\varepsilon)
\;=\; 1-\varepsilon' .
\]
On the rounds of $R'_d$ the emitted marks are the type's own marks, and the revealing value is
drawn from the $\kappa'$ conditional law computed in the proof of Lemma~\ref{lem:g1}. By that
lemma's factorisation across rounds, the emitted history has exactly the $\kappa'$ law given the
type. So $\Lambda$ carries the $\kappa$ experiment to the $\kappa'$ experiment and the second is
a garbling of the first.

For the final claim, run $\Lambda$ on the $\kappa$ history to build a joint law of the type, the
$\kappa$ history and the $\kappa'$ history. Because $\Lambda$ reads only the history, those three
form a Markov chain in that order, so
\[
\Prb\bigl(\theta \in \cdot \bigm| X^{\kappa'}_{0:d}\bigr)
\;=\; \Eop\Bigl[\,\Prb\bigl(\theta \in \cdot \bigm| X^{\kappa}_{0:d}\bigr)
\Bigm| X^{\kappa'}_{0:d}\Bigr].
\]
The $\kappa'$ posterior is therefore a conditional expectation of the $\kappa$ posterior, so the
law of the first is a mean preserving contraction of the law of the second, and when
$\Omega(\tau,T) < 1$, both have mean $\rho_P^{\tau,T}$.
\end{proof}

\begin{proof}[Proof of Lemma~\ref{lem:g3}]
Part~(a). By Lemma~\ref{lem:g1}(a) the revealed set $R_H$ has the product Bernoulli law
$\Prb(R_H = S) = (1-\varepsilon)^{|S|}\varepsilon^{\,H+1-|S|}$, and that law is the same under
$\rho$ and under $\rho_P^{\tau,T}$, because $\{D=0\}$ is a type event and $R_H$ is independent of
the type. By Lemma~\ref{lem:g1}(c) the pooled posterior on $\{R_H = S\}$ is
$\rho_P^{\tau,T}(\cdot \mid (m_e)_{e \in S})$. Conditioning on $R_H$ and applying the tower
property,
\[
\Eop\bigl[h(\nu_{0:H}) \mid D=0\bigr]
\;=\; \sum_{S} \Prb(R_H = S)\;
\Eop_{\rho_P^{\tau,T}}\Bigl[h\bigl(\rho_P^{\tau,T}(\cdot\mid (m_e)_{e\in S})\bigr)\Bigr]
\;=\; \sum_S \Prb(R_H = S)\,\Grev_{\tau,T}(S),
\]
and multiplying by $\Delta_m$ gives \eqref{eq:g-rep}. Nothing in the argument used the depth or
the particular cell event, so the same identity holds with $H$ replaced by any $d \le H$ and
$\{D=0\}$ replaced by any cell event. The numbers $\Grev_{\tau,T}(S)$ are built from
$\rho_P^{\tau,T}$ and the mark paths only, and neither depends on $\kappa$.

Part~(b). Write $q := 1-\varepsilon$ and $F(q) := \sum_S q^{|S|}(1-q)^{H+1-|S|}\Grev(S)$, the
expectation of $\Grev(R_H)$ under the product Bernoulli measure with success probability $q$ on
each of the $H+1$ rounds. Read $F$ as a function of $H+1$ separate success probabilities, one
per round. It is affine in each of them, because the rounds enter the product measure
independently, and its partial derivative in the $d$-th is
$\Eop[\Grev(R_{-d}\cup\{d\})] - \Eop[\Grev(R_{-d})]$, with $R_{-d}$ the random set on the other
rounds. Setting every coordinate to $q$ and summing the $H+1$ partial derivatives,
\[
F'(q) \;=\; \sum_{d=0}^{H}
\Eop_{q}\Bigl[\,\Grev\bigl(R_{-d}\cup\{d\}\bigr) - \Grev\bigl(R_{-d}\bigr)\Bigr],
\]
Expanding the expectation over $R_{-d}$ and collecting terms by $|S| = k$,
\[
F'(q) \;=\; \sum_{d=0}^{H}\ \sum_{S \subseteq \{0,\ldots,H\}\setminus\{d\}}
q^{|S|}(1-q)^{H-|S|}\bigl[\Grev(S\cup\{d\})-\Grev(S)\bigr]
\;=\; \sum_{k=0}^{H} q^{k}(1-q)^{H-k} c_k(\tau,T).
\]
Since $M_P = \Delta_m F(q)$ and $q = 1-\kappa/2$ gives $dq/d\kappa = -1/2$, the chain rule yields
\eqref{eq:g-deriv} with $1-\varepsilon = q$ and $\varepsilon = 1-q$.

Part~(c). Fix $S$ and $d \notin S$, and write $\nu_S$ and $\nu_{S\cup\{d\}}$ for the two
posteriors. The record on $S$ is coarser than the record on $S\cup\{d\}$, so
$\nu_S = \Eop[\nu_{S\cup\{d\}} \mid (m_e)_{e\in S}]$. Both $\pi$ and $\hat v$ are linear in the
posterior, so the pair $(\pi,\hat v)$ evaluated at these two posteriors is an $\mathbb R^2$ valued
martingale over the two records, and its values lie in the convex hull named in the statement. If
$\mathsf h$ is concave there, the conditional Jensen inequality in $\mathbb R^2$ gives
\[
\Eop\bigl[\mathsf h(\pi_{S\cup\{d\}},\hat v_{S\cup\{d\}}) \bigm| (m_e)_{e\in S}\bigr]
\;\le\; \mathsf h(\pi_S, \hat v_S)
\quad\text{almost surely,}
\]
and taking expectations, $\Grev(S\cup\{d\}) \le \Grev(S)$. Every increment in \eqref{eq:g-ck} is
then at most zero, so $c_k \le 0$ for every $k$. The weights $(1-\varepsilon)^k\varepsilon^{H-k}$
are strictly positive on $\kappa \in (0,1)$ and $\Delta_m > 0$, so \eqref{eq:g-deriv} gives
$\partial_\kappa M_P \ge 0$ there. Reversing the curvature reverses each of the three steps. The
depth played no role, so the statement holds at every $d \le H$.
\end{proof}

\begin{proof}[Proof of Theorem~\ref{thm:g-threshold}]
Part~(A). At fixed policies the disclosure indicator is a function of the type alone: the plan
fixes the stake path, the path fixes the crossing date, and the filing lands by the control date
exactly when the crossing date is at most $H-T$. Liquidity enters only the noise law, which
touches order flow and not the stake path, so $\Omega$ does not move with $\kappa$. The flagged
average $M_F$ does not move with $\kappa$ by the flagged invariance assumed in the statement.
Differentiating $\Dact = \Omega M_F + (1-\Omega)M_P$,
\[
\partial_\kappa \Dact
= (\partial_\kappa \Omega)(M_F - M_P) + \Omega\,\partial_\kappa M_F
  + (1-\Omega)\,\partial_\kappa M_P
= (1-\Omega)\,\partial_\kappa M_P ,
\]
the first two terms vanishing term by term rather than by cancellation. Differentiability of
$M_P$ in $\kappa$ is not a hypothesis here: Lemma~\ref{lem:g3}(a) exhibits $M_P$ as a polynomial
in $\kappa$. Taking absolute values and using $1-\Omega \ge 0$ gives
$\mathcal S = (1-\Omega)\mathcal S_P$.

Part~(B). Let a history be flagged at $\tau$, so $a = 1$ and $c(\tau) + T \le H$, where
$c(\tau) = \inf\{d : B(s,d) \ge \tau\}$. Since $\tau' < \tau$, every date at which the stake
reaches $\tau$ is a date at which it reaches $\tau'$, so $c(\tau') \le c(\tau)$ and therefore
$c(\tau') + T \le H$. The engagement flag is unchanged, so the history is flagged at $\tau'$ as
well, and $\mathcal C_F(\tau,T)\subseteq\mathcal C_F(\tau',T)$. Disclosure implies engagement, so
every history in the difference carries $a=1$ and is generated by a Voice plan. Taking masses,
$\Omega(\tau',T)\ge\Omega(\tau,T)$, and with $\Omega(\tau',T) < 1$,
\[
W_\tau = \frac{1-\Omega(\tau',T)}{1-\Omega(\tau,T)} \in [0,1],
\]
with $W_\tau = 1$ exactly when the two flagged sets have the same mass, that is when no history
is reclassified (almost surely). The argument used no property of the kernel.

Part~(C). Both thresholds are evaluated at the same $\kappa$ and at the same policies, so
Lemma~\ref{lem:g3}(b) applies at each and gives
\[
\mathcal S_P(\tau,T) = \frac{\Delta_m}{2}\bigl|\Wrev_{\tau,T}(\kappa)\bigr|,
\qquad
\mathcal S_P(\tau',T) = \frac{\Delta_m}{2}\bigl|\Wrev_{\tau',T}(\kappa)\bigr| .
\]
Since $\mathcal S_P(\tau,T) > 0$ the ratio $C_\tau$ is defined, and
Condition~\ref{cond:D} states exactly that its numerator is at most its denominator, so
$C_\tau \le 1$. Part~(A) applied at each threshold gives
$\mathcal S(\tau',T)/\mathcal S(\tau,T) = W_\tau C_\tau$, and part~(B) puts $W_\tau$ in $[0,1]$.
The product of a number in $[0,1]$ with a number in $[0,1]$ is at most one, which is
\eqref{eq:g-dial}, with equality exactly when both factors are one.
\end{proof}

\begin{proof}[Proof of Remark~\ref{rem:g-Dstar}]
Write $w_k = (1-\varepsilon)^k\varepsilon^{H-k} > 0$. By the triangle inequality and
\eqref{eq:g-Dstar},
\[
\bigl|\Wrev_{\tau',T}(\kappa)\bigr| \;\le\; \sum_{k} w_k\,|c_k(\tau',T)|
\;\le\; \sum_{k} w_k\,|c_k(\tau,T)| .
\]
If the numbers $c_k(\tau,T)$ share one sign then the last sum equals
$\bigl|\sum_k w_k c_k(\tau,T)\bigr| = |\Wrev_{\tau,T}(\kappa)|$. The bound holds at every
$\kappa \in (0,1)$ because the hypothesis does not involve $\kappa$.
\end{proof}

\begin{remark}[Verification of Condition~\ref{cond:D}]
\label{rem:g-grid}
Condition~\ref{cond:D} is checked, and the coefficients $c_k$ of \eqref{eq:g-ck} are reported, at
order size two on the calibration grid: the frozen baseline policy, the two windows
$T \in \{5,10\}$, the five threshold quantiles $0.1$, $0.3$, $0.5$, $0.7$, $0.9$ of the
terminal stake distribution under the benchmark policy, taken in adjacent pairs, and the $71$
node liquidity grid $\kappa = 0.15, 0.16, \ldots, 0.85$. The record is
\texttt{numerical\_v4/checks/t2\_threshold\_revelation\_check.json}. Its status on that grid is
NUMERICAL.
\end{remark}


% =============================================================================
% proofs/07_detection.tex
% Detection, silence, and bidder entry at fixed policies.
% Fragment, to be \input by appendix.tex.
% =============================================================================

\section{Detection: the entry identity, the tape, who gets caught and silence}
\label{sec:detection}

This section connects the engagement premium to bidder entry and separates the two ways public
information identifies a Voice type. The filing acts through the disclosure cell. The tape acts
through flow values that reveal a building purchase.
At any date $d$, write $D^d$ for the indicator that the filing has landed by $d$, and write
$\mathcal C_F^d=\{D^d=1\}$ and $\mathcal C_P^d=\{D^d=0\}$.

\csname unlabelledtrue\endcsname
\begin{lemma}[Entry identity]
\label{lem:entry-identity}
Let $\mathcal I_H$ be the public information set at the control node, let
$\pi(\mathcal I_H)=\Prb(a=1\mid\mathcal I_H)$, and let
$p(\mathcal I_H)$ be the bidder-entry probability. The premium kernel is
$h(\mathcal I_H)=\pi(\mathcal I_H)p(\mathcal I_H)$. Assume:
\begin{enumerate}[label=(E-\arabic*),leftmargin=3.2em,itemsep=3pt]
\item \emph{Public measurability.} The price and the entry probability are functions of the public
  information set, and the engagement indicator is a function of the type. Thus
  $P=P(\mathcal I_H)$ and $p=p(\mathcal I_H)$ are $\mathcal I_H$-measurable, while
  $a=a(\theta)$. The first assertion is the pinned-kernel content of (S8); the last follows from
  the finite menu and cutoff selection map in (S2) and (S3).
\item \emph{The Voice event and the public cells.} The menu assigns $a=1$ exactly to Voice, and
  $a=0$ to Exit and Hold. The cell indicator $D$ is public by (S5) and (S6), so
  $\mathcal C_F=\{D=1\}$ and $\mathcal C_P=\{D=0\}$ belong to $\mathcal I_H$.
\item \emph{Finite premium and fixed policies.} The premium satisfies (S9), liquidity enters only
  through the noise law as in (S10), and the menu, cutoff vector, and execution policies are fixed
  as in (S11).
\end{enumerate}
Write $\omega_V=\Prb(a=1)$ for the Voice mass. If $\omega_V>0$, define
$e_V=\Eop[p(\mathcal I_H)\mid a=1]$. Then
\begin{equation}
\label{eq:entry-identity}
  \Dact=\Delta_m\,\omega_V e_V .
\end{equation}
For either cell $\mathcal C\in\{\mathcal C_F,\mathcal C_P\}$ with
$\Prb(\mathcal C)>0$ and $\Prb(a=1,\mathcal C)>0$, put
\[
  \omega_{V\mid\mathcal C}=\Prb(a=1\mid\mathcal C),
  \qquad
  e_{V\mid\mathcal C}=\Eop[p(\mathcal I_H)\mid a=1,\mathcal C].
\]
Its cell average obeys the same identity,
\begin{equation}
\label{eq:entry-identity-cell}
  M_{\mathcal C}
  :=\Delta_m\Eop[\pi(\mathcal I_H)p(\mathcal I_H)\mid\mathcal C]
  =\Delta_m\,\omega_{V\mid\mathcal C}e_{V\mid\mathcal C}.
\end{equation}
If the Voice mass is zero, the corresponding premium is zero and the conditional entry
probability is left undefined. At fixed policies, $\omega_V$ is independent of both $\kappa$ and
the disclosure rule.
\end{lemma}
\csname unlabelledfalse\endcsname

\begin{proof}
By (E-1), $p(\mathcal I_H)$ is $\mathcal I_H$-measurable, and (S8) makes it bounded. The defining
property of conditional expectation therefore gives
\[
  \Eop[\pi(\mathcal I_H)p(\mathcal I_H)]
  =\Eop\!\left[\Eop[a\mid\mathcal I_H]p(\mathcal I_H)\right]
  =\Eop[a\,p(\mathcal I_H)].
\]
When $\omega_V>0$, conditioning the last expectation on the event $\{a=1\}$ gives
$\Eop[a p]=\omega_V\Eop[p\mid a=1]=\omega_Ve_V$. Multiplication by $\Delta_m$ and (S9) yield
\eqref{eq:entry-identity}. If $\omega_V=0$, then $a=0$ almost surely, so both
$\Eop[a p]$ and the engagement premium vanish.

Now fix a cell $\mathcal C$ of positive mass. Its indicator is $\mathcal I_H$-measurable by
(E-2). Applying the same calculation after multiplication by $\ind_{\mathcal C}$ gives
\[
  \Eop[\pi p\,\ind_{\mathcal C}]
  =\Eop[a p\,\ind_{\mathcal C}].
\]
Divide by $\Prb(\mathcal C)$. If $\Prb(a=1,\mathcal C)>0$, conditioning the right side first on
$\{a=1\}\cap\mathcal C$ gives \eqref{eq:entry-identity-cell}. If that joint mass is zero, the
right side is zero, hence so is the cell premium.

It remains to check the invariance of the unconditional Voice mass. By (S3) and (S7),
$a=a_{j(s)}$. Condition (S11) holds the menu and the cutoff map $j(\cdot)$ fixed across rules and
in $\kappa$, so neither the threshold margin nor the window margin enters this indicator.
Condition (S10) puts $\kappa$ only in the law of the noise marks, while the law of the signal is
unchanged by (S1). The law of $a_{j(s)}$, and hence $\omega_V$, is therefore common to every rule
and every $\kappa$ in a fixed-policy comparison.
\end{proof}

\csname unlabelledtrue\endcsname
\begin{lemma}[Detection]
\label{lem:detection}
Fix a date $d\in\{0,\ldots,H\}$ and a pooled history of positive probability. Assume:
\begin{enumerate}[label=(D-\arabic*),leftmargin=3.2em,itemsep=3pt]
\item \emph{Voice marks.} The order mark is $g_e(\theta)\in\{0,1\}$ and strategic order flow is
  $q_e(\theta)=2\bar z\,g_e(\theta)$. Only Voice plans carry positive marks. Exit and Hold have
  $g_e=0$ at every date, and a Voice building round has $g_e=1$.
\item \emph{Common noise channel.} Condition (S12) holds. In particular,
  $X_e=q_e+z_e$, the noise marks are independent across dates and independent of the type, and
  $z_e$ equals $-\bar z$, $0$, or $+\bar z$ with probabilities
  $\kappa/2$, $1-\kappa$, and $\kappa/2$, where $\kappa\in[0,1]$.
\item \emph{Public observation.} Conditions (S6), (S7), (S8), and (S13) hold, so the pooled flow
  history is public, its posterior is the pinned conditional probability, and the date-$d$ pooled
  cell is a type event.
\end{enumerate}
The flow values $2\bar z$ and $3\bar z$ reveal a building mark. On any positive-probability
pooled history containing either value, the engagement posterior is one. The values $-\bar z$
and $0$ reveal a zero mark whenever those values have positive probability. The value $+\bar z$
is the only ambiguous value: in the round-level channel support at $\kappa>0$, it can arise from a
building purchase with a noise sale or from a zero mark with a noise purchase. At $\kappa=0$ the
ambiguous value is not observed.

For a Voice type, let
\[
  N_d(\theta)=\sum_{e=0}^{d}g_e(\theta)
\]
be the number of building rounds through that date. The probability that no building mark has
been revealed on the tape is
\begin{equation}
\label{eq:detection-silence}
  \Prb(\text{no building mark revealed through }d\mid\theta)
  =\left(\frac{\kappa}{2}\right)^{N_d(\theta)}.
\end{equation}
The right side is the product over building rounds, with an empty product equal to one.
If the type remains pooled at date $d$, this event is a silent history and the same formula is its
silent probability.
\end{lemma}
\csname unlabelledfalse\endcsname

\begin{proof}
If $g_e=0$, then $q_e=0$, so the three possible flow values are
$-\bar z$, $0$, and $+\bar z$. If $g_e=1$, then $q_e=2\bar z$, so the three possible values are
$+\bar z$, $2\bar z$, and $3\bar z$. The two supports meet only at $+\bar z$. This proves the
stated classification.

The event $\{X_e\in\{2\bar z,3\bar z\}\}$ implies $g_e=1$. By (D-1), it therefore implies that
the selected plan is Voice and $a=1$. On a pooled history of positive probability that contains
such a value, Bayes' rule assigns probability one to $a=1$. Thus
$\pi(\mathcal I_d)=1$ on that history.

Fix a Voice type. On each of its $N_d(\theta)$ building rounds, the only flow value
that does not reveal the positive mark is $+\bar z$. That value occurs on a building round
exactly when $z_e=-\bar z$, an event of probability $\kappa/2$. Zero-mark rounds cannot produce
$2\bar z$ or $3\bar z$ and hence cannot reveal a building mark. Independence across dates in
(S12) now gives the product in \eqref{eq:detection-silence}. The event $D^d=0$ is a function of
the type under (S7), so restricting the formula to pooled types does not change the noise product
law.
\end{proof}

\csname unlabelledtrue\endcsname
\begin{lemma}[Upper set]
\label{lem:upper-set}
Fix a threshold margin $\tau$, a window margin $T$, and a date
$d\in\{0,\ldots,H\}$. Assume:
\begin{enumerate}[label=(U-\arabic*),leftmargin=3.2em,itemsep=3pt]
\item \emph{Monotone target.} The Voice target $b^*(s)\ge b_0$ is weakly increasing in the
  signal.
\item \emph{Monotone accumulation length.} The accumulation length
  $n(s)\in\{1,\ldots,H+1\}$ is weakly decreasing in the signal.
\item \emph{Ordered menu form.} Voice is selected on an upper interval of the signal line, and
  its stake and mark paths are
  \[
    B_V(s,e)=b_0+\bigl(b^*(s)-b_0\bigr)
      \min\!\left\{1,\frac{e+1}{n(s)}\right\},
    \qquad
    g_e(s)=\ind\{e<n(s)\}.
  \]
\item \emph{Timing and tape.} Conditions (S4), (S5), and (S7) hold, together with (D-1) and the
  common noise channel (S12).
\end{enumerate}
Then the date-$d$ flagged signal set
\begin{equation}
\label{eq:upper-flagged-set}
  \mathcal F_d
  =\{s:j(s)=\mathrm{Voice},\ c(s;\tau)+T\le d\}
\end{equation}
is an upper set. If $s_2\ge s_1$ are Voice signals, then
\[
  c(s_2;\tau)\le c(s_1;\tau),
  \qquad
  N_d(s_2)\le N_d(s_1),
  \qquad
  N_d(s)=\min\{n(s),d+1\}.
\]
Consequently, for prior-almost every pair of Voice types $\theta_F$ and $\theta_P$ with
$s(\theta_F)\in\mathcal F_d$ and
$s(\theta_P)\notin\mathcal F_d$,
\begin{equation}
\label{eq:upper-silence-order}
  \Prb(\text{no building mark revealed through }d\mid\theta_F)
  \ge
  \Prb(\text{no building mark revealed through }d\mid\theta_P).
\end{equation}
The inequality is strict when $\kappa>0$ and $N_d(\theta_F)<N_d(\theta_P)$. Thus the rule flags from
the Voice types with weakly fewer building rounds, which occupy the weakly highest-probability end
among Voice types for escaping tape identification.
\end{lemma}
\csname unlabelledfalse\endcsname

\begin{proof}
Take Voice signals $s_2\ge s_1$. By (U-1),
$b^*(s_2)-b_0\ge b^*(s_1)-b_0\ge0$. By (U-2),
$n(s_2)\le n(s_1)$, and therefore, at every date $e$,
\[
  \min\!\left\{1,\frac{e+1}{n(s_2)}\right\}
  \ge
  \min\!\left\{1,\frac{e+1}{n(s_1)}\right\}.
\]
Both factors in the product defining the increment above $b_0$ are nonnegative and weakly
larger at $s_2$. The path formula in (U-3) gives
$B_V(s_2,e)\ge B_V(s_1,e)$ for every $e$.

Under (S5), $c(s;\tau)$ is the first date at which this path reaches $\tau$. If the lower signal
has crossed by a date, the higher signal has crossed by that date as well. Hence
$c(s_2;\tau)\le c(s_1;\tau)$. If $s_1\in\mathcal F_d$, then (U-3) keeps every higher signal in
Voice and the crossing inequality gives
$c(s_2;\tau)+T\le c(s_1;\tau)+T\le d$. Thus $s_2\in\mathcal F_d$, proving that
$\mathcal F_d$ is an upper set. The path is fixed before flow is observed by (S7), so this
ordering is not altered by realised noise or prices.

The mark formula in (U-3) makes the mark equal to one at each date
$e<\min\{n(s),d+1\}$ and none after it, so
$N_d(s)=\min\{n(s),d+1\}$. This count is weakly decreasing in $s$ by (U-2). Since
$\mathcal F_d$ is upper, a flagged signal cannot lie below an unflagged Voice signal. Thus, for the
types in the statement, $N_d(\theta_F)\le N_d(\theta_P)$. Lemma~\ref{lem:detection} gives the two
no-revelation probabilities as $(\kappa/2)^{N_d(\theta_F)}$ and
$(\kappa/2)^{N_d(\theta_P)}$. The base lies in $[0,1/2]$, so the smaller
exponent gives the weakly larger probability. It gives a strict inequality when the base is
positive and the exponents differ.
\end{proof}

\csname unlabelledtrue\endcsname
\begin{lemma}[Silence]
\label{lem:silence}
Fix a date $d\in\{0,\ldots,H\}$ and compare rules at common fixed policies and common $\kappa$.
Either $r_+=(\tau',T)$ and $r_-=(\tau,T)$ with $b_0<\tau'<\tau$, or
$r_+=(\tau,T')$ and $r_-=(\tau,T)$ with $T'<T$. Let
\[
  A=\mathcal C_{P,r_+}^{d},
  \qquad
  R=\mathcal C_{P,r_-}^{d}\setminus\mathcal C_{P,r_+}^{d}.
\]
Assume:
\begin{enumerate}[label=(L-\arabic*),leftmargin=3.2em,itemsep=3pt]
\item \emph{Nested upper removal.} Conditions (S4), (S5), (S7), and (S11), together with
  (U-1) to (U-3), imply
  $\mathcal C_{P,r_+}^{d}\subseteq\mathcal C_{P,r_-}^{d}$. The set $R$ has positive prior mass,
  and $A$ and $R$ are $\sigma(s)$-measurable at fixed policies by Step~4 of
  Lemma~\ref{lem:partition}. The set $R$ lies in $\{a=1\}$ and is upper relative to $A$: for
  prior-almost every
  $\theta_R\in R$ and $\theta_A\in A$, $s(\theta_R)\ge s(\theta_A)$.
\item \emph{Common pooled likelihood.} Conditions (S12) and (S13) hold. On $A$, the two rules
  have the same type-conditional law of the displayed pooled flow history.
\item \emph{Monotone signal information.} The conditional mean
  $\mu_v(s)=\Eop[v\mid s]$ exists and is weakly increasing in $s$, and
  $v$ is conditionally independent of pooled flow given $s$. In the Gaussian signal model,
  $\mu_v(s)=\mu_v+\beta(s-\mu_v)$ with $\beta>0$, so this ordering is strict.
\item \emph{Common positive-probability history.} The displayed pooled flow history
  $x=(x_0,\ldots,x_d)$, with all public filing coordinates equal to zero, has positive
  probability under both rules.
\end{enumerate}
Let
\[
  \pi_r(x)=\Prb(a=1\mid X_{0:d}=x,D_r^d=0),
  \qquad
  \widehat v_r(x)=\Eop[v\mid X_{0:d}=x,D_r^d=0].
\]
Then
\begin{equation}
\label{eq:silence-posterior-order}
  \pi_{r_+}(x)\le\pi_{r_-}(x),
  \qquad
  \widehat v_{r_+}(x)\le\widehat v_{r_-}(x).
\end{equation}
The first inequality is strict exactly when the looser posterior assigns positive probability to
$R$ and $\pi_{r_+}(x)<1$. The second is strict when the looser posterior assigns positive
probability to $R$ and the posterior mean of $\mu_v(s)$ on $R$ is strictly larger than its
posterior mean on $A$. In particular, strict monotone signal information and strictly separated
posterior signal supports give strictness when the looser posterior assigns positive probability
to $R$. The conclusions apply to every silent history that satisfies (L-4), as well as to every
other common pooled history of positive probability.
\end{lemma}
\csname unlabelledfalse\endcsname

\begin{proof}
Nestedness partitions the looser pooled type set into $A$ and $R$. Condition (S11) gives both rules
the same prior over types. By (L-2), the likelihood of
the common flow vector $x$ given a type is the same under both rules. Conditioning the looser
posterior further on $A$ therefore gives the tighter posterior. Put
\[
  \lambda_x=\Prb(R\mid X_{0:d}=x,D_{r_-}^d=0).
\]
Assumption (L-4) gives positive posterior mass to $A$, so $0\le\lambda_x<1$. If
$\lambda_x=0$, the two posteriors coincide and both inequalities in
\eqref{eq:silence-posterior-order} are equalities. Suppose now that $\lambda_x>0$.

Every type in $R$ is engaged by (L-1). The looser engagement posterior is therefore the mixture
\[
  \pi_{r_-}(x)
  =(1-\lambda_x)\pi_{r_+}(x)+\lambda_x.
\]
Since $\pi_{r_+}(x)\le1$, this proves the first inequality. The displayed mixture also shows that
it is strict exactly under the conditions stated.

Let
\[
  \widehat v_R(x)=\Eop[v\mid X_{0:d}=x,\theta\in R].
\]
The same posterior decomposition gives
\begin{equation}
\label{eq:silence-v-mixture}
  \widehat v_{r_-}(x)
  =(1-\lambda_x)\widehat v_{r_+}(x)+\lambda_x\widehat v_R(x).
\end{equation}
By (L-3), conditioning on the flow adds no information about $v$ after conditioning on $s$.
Thus the two component means in \eqref{eq:silence-v-mixture} are averages of $\mu_v(s)$ under
their respective posterior signal laws. Assumption (L-1) puts every signal in $R$ weakly above
every signal in $A$, and (L-3) makes $\mu_v$ weakly increasing. Hence
$\widehat v_R(x)\ge\widehat v_{r_+}(x)$. Substitution into
\eqref{eq:silence-v-mixture} proves the second inequality. It is strict exactly when
$\lambda_x>0$ and the component means differ, with the stated signal conditions sufficient for
that difference.
\end{proof}

\csname unlabelledtrue\endcsname
\begin{lemma}[Entry against the undetected]
\label{lem:entry-undetected}
Let $\Pi(\widehat v,\pi)$ be the unique pricing root of
Lemma~\ref{lem:pricing-root}, and define the entry probability at that root by
$e(\widehat v,\pi)=p(\Pi(\widehat v,\pi),\pi)$. Fix an interval
$\mathcal V\subset\mathbb R$ of fundamental posterior means. Assume:
\begin{enumerate}[label=(I-\arabic*),leftmargin=3.2em,itemsep=3pt]
\item \emph{Unique pinned root.} The hypotheses of Lemma~\ref{lem:pricing-root} hold, including
  $m_0\ge0$, $\Delta_m>0$, $\Delta_V\ge0$, $\sigma_\xi>0$, and uniqueness of the finite fixed
  point. The pinned pricing rule is the one in (S8).
\item \emph{Entry derivatives.} For every finite $P$ and $\pi\in[0,1]$, entry has the normal form
  \[
    p(P,\pi)=1-\Phi\!\left(
      \frac{P+K+m_0+\pi\Delta_m-\bar S}{\sigma_\xi}\right),
  \]
  so it is continuously differentiable with
  $\partial_Pp<0$ and $\partial_\pi p<0$.
\item \emph{Added engagement-price condition.} At every root
  $P=\Pi(\widehat v,\pi)$ with $(\widehat v,\pi)\in\mathcal V\times[0,1]$, write
  $m(\pi)=m_0+\pi\Delta_m$ and $o(P,\pi)=p(P,\pi)/(1-p(P,\pi))$. Assume
  \begin{equation}
  \label{eq:entry-price-condition}
    \Delta_V+\Delta_m o(P,\pi)+m(\pi)\,\partial_\pi o(P,\pi)>0.
  \end{equation}
  Equivalently, with
  $t=(P+K+m(\pi)-\bar S)/\sigma_\xi$ and standard normal density $\phi$,
  \[
    \Delta_V+\Delta_m o(P,\pi)
    >
    \frac{m(\pi)\Delta_m\phi(t)}{\sigma_\xi\Phi(t)^2}.
  \]
\item \emph{Added history-comparison condition.} For the history comparison, the silent pooled
  history has engagement posterior $\pi_U\le1$. A tape-detected pooled history with the same fundamental
  posterior mean has engagement posterior one. The comparison is nondegenerate when $\pi_U<1$.
\end{enumerate}
Then $\Pi(\widehat v,\pi)$ is strictly increasing in both $\widehat v$ and $\pi$ on the stated
domain, while $e(\widehat v,\pi)$ is strictly decreasing in both. Hence, holding the fundamental
posterior mean fixed,
\begin{equation}
\label{eq:entry-undetected-gap}
  e(\widehat v,\pi_U)\ge e(\widehat v,1),
  \qquad \widehat v\in\mathcal V.
\end{equation}
The inequality is strict when $\pi_U<1$ and is an equality when $\pi_U=1$.
\end{lemma}
\csname unlabelledfalse\endcsname

\begin{proof}
The pricing fixed point can be written as the root of
\begin{equation}
\label{eq:entry-root-residual}
  \mathcal R(P;\widehat v,\pi)
  =P-\widehat v-\pi\Delta_V-m(\pi)o(P,\pi)=0.
\end{equation}
By (I-2), $p\in(0,1)$, $\partial_Pp<0$, and the odds map is strictly increasing in $p$.
Therefore $\partial_Po<0$. Since $m(\pi)\ge0$,
\begin{equation}
\label{eq:entry-root-slope}
  \partial_P\mathcal R=1-m(\pi)\partial_Po>0.
\end{equation}
Lemma~\ref{lem:pricing-root} supplies the unique root. Equation
\eqref{eq:entry-root-slope} lets the implicit function theorem differentiate that root.
Because $\partial_{\widehat v}\mathcal R=-1$,
\[
  \partial_{\widehat v}\Pi
  =-\frac{\partial_{\widehat v}\mathcal R}{\partial_P\mathcal R}
  =\frac{1}{\partial_P\mathcal R}>0.
\]
For the engagement posterior,
\[
  \partial_\pi \mathcal R
  =-\Delta_V-\Delta_mo-m(\pi)\partial_\pi o.
\]
Condition \eqref{eq:entry-price-condition} says exactly that this derivative is negative. Hence
\[
  \partial_\pi\Pi
  =-\frac{\partial_\pi \mathcal R}{\partial_P\mathcal R}>0.
\]
Continuity of the unique root, supplied by Lemma~\ref{lem:pricing-root}, extends these monotone
comparisons to the endpoints of the posterior interval.

Apply the chain rule to entry at the root:
\[
  \partial_{\widehat v}e
  =(\partial_Pp)(\partial_{\widehat v}\Pi)<0,
  \qquad
  \partial_\pi e
  =(\partial_Pp)(\partial_\pi\Pi)+\partial_\pi p<0.
\]
Both signs use all the derivative restrictions in (I-2); the second also uses (I-3). Under (I-4),
decrease in $\pi$ at fixed $\widehat v$ gives \eqref{eq:entry-undetected-gap}, strictly when
$\pi_U<1$. When $\pi_U=1$, both arguments of $e$ are the same.
\end{proof}


\section{Who gets caught: nested cuts under both dials}
\label{sec:app-caught}

% proofs/03_caught.tex
% Who gets caught: signing the composition ratio of the window margin.
% Fragment, to be \input by appendix.tex.
%
% appendix.tex must supply: amsmath, amssymb, amsthm (the proof environment), enumitem,
%   \newtheorem{corollary}{Corollary}, and the macros \Eop, \Prb, \ind.
% Cross-references to the partition result, the clock equivalence and the clock
%   theorem are written in words here.
% Notation notes: (a) B below is the set of newly caught histories. It
%   is unrelated to the stake symbols B_j(s,t) and B^F. (b) The caught share is
%   written \varphi, not \beta, because \beta is the signal weight
%   sigma_v^2/(sigma_v^2 + sigma_eps^2) in the model section. Rename either if the
%   collision reads badly.

\subsection{Who gets caught}
\label{sec:caught}

The clock theorem makes window attenuation equivalent to $W_T C_T \le 1$ and signs only the weight
leg. This section signs the composition leg. It writes the longer clock's pooled sensitivity as a
mass-weighted average of what the shorter clock leaves in the pool and what it takes out, and reads
both $C_T \le 1$ and $W_T C_T \le 1$ off that average.

\paragraph{Setup.}
Fix the threshold margin $\tau$ and the plan and cutoff policies, and compare two window margins
$T' < T$ at a common $\kappa$. Write
\begin{equation}
\label{eq:caught-sets}
A \;=\; \mathcal{C}_P(\tau,T),
\qquad
B \;=\; \mathcal{C}_F(\tau,T')\setminus\mathcal{C}_F(\tau,T),
\end{equation}
so $A$ is the pooled cell under the longer clock and $B$ collects the histories the shorter clock
newly flags. Write
\begin{equation}
\label{eq:caught-phi}
\varphi \;=\; \frac{\Prb(B)}{\Prb(A)}
\end{equation}
for the share of the longer clock's pooled cell that the shorter clock catches.

Each rule pins its own premium kernel. The control-node information set carries the coordinate
``the filing has landed by date $d$'', so the pooled history at window margin $T$ and the pooled
history at window margin $T'$ are different observables, and the engagement posterior
$\pi = \Prb(a = 1 \mid \cdot)$ and the price are taken with respect to whichever of the two the
rule supplies. Write $h^{(T)}$ and $h^{(T')}$ for the kernel $h = \pi p$ evaluated at the
control-node information set of the rule $(\tau,T)$ and of the rule $(\tau,T')$. The two pooled
averages of the clock theorem are then
\begin{equation}
\label{eq:caught-MP}
M_P(\tau,T) \;=\; \Delta_m\,\Eop\bigl[h^{(T)} \bigm| A\bigr],
\qquad
M_P(\tau,T') \;=\; \Delta_m\,\Eop\bigl[h^{(T')} \bigm| A\setminus B\bigr],
\end{equation}
the second reading of the pooled cell at the shorter clock being part~(i) below. Set
\begin{equation}
\label{eq:caught-sens}
s_A \;=\; \partial_\kappa\Eop\bigl[h^{(T)} \bigm| A\bigr],
\qquad
s_{A\setminus B} \;=\; \partial_\kappa\Eop\bigl[h^{(T')} \bigm| A\setminus B\bigr],
\end{equation}
so that, with $\Delta_m > 0$ finite,
\begin{equation}
\label{eq:caught-SP}
\mathcal{S}_P(\tau,T) \;=\; \Delta_m\,|s_A|,
\qquad
\mathcal{S}_P(\tau,T') \;=\; \Delta_m\,|s_{A\setminus B}|,
\qquad
C_T \;=\; \frac{|s_{A\setminus B}|}{|s_A|} ,
\end{equation}
by the definitions of the two pooled sensitivities and no more.

Two further quantities separate what the catch removes from how the surviving pool is re-read:
\begin{equation}
\label{eq:caught-raw}
\tilde s_B \;=\; \partial_\kappa\Eop\bigl[h^{(T)} \bigm| B\bigr],
\qquad
\delta \;=\; \partial_\kappa\Eop\bigl[h^{(T')} - h^{(T)} \bigm| A\setminus B\bigr] .
\end{equation}
Here $\tilde s_B$ is the \emph{caught-only leg}, the noise sensitivity of the newly caught
histories priced as the longer clock priced them, and $\delta$ is the \emph{re-pricing term}: the pooled posterior conditions
on the pool the rule leaves behind, so the histories that stay pooled are read against a smaller
pool once the clock shortens, and $\delta$ is the rate at which that re-reading moves their noise
sensitivity. Define the \emph{net cut leg} by
\begin{equation}
\label{eq:caught-sB}
s_B \;:=\; \frac{1}{\Prb(B)}\,
\partial_\kappa\Bigl(
\Eop\bigl[h^{(T)}\ind_A\bigr] - \Eop\bigl[h^{(T')}\ind_{A\setminus B}\bigr]
\Bigr) .
\end{equation}
The pooled cell contributes to the expected engagement premium an amount
\[
\begin{aligned}
\Delta_m\,\Eop\bigl[h^{(T)}\ind_A\bigr]
  \;&=\; \bigl(1-\Omega(\tau,T)\bigr)\,M_P(\tau,T)
  &&\text{at the longer clock,} \\[3pt]
\Delta_m\,\Eop\bigl[h^{(T')}\ind_{A\setminus B}\bigr]
  \;&=\; \bigl(1-\Omega(\tau,T')\bigr)\,M_P(\tau,T')
  &&\text{at the shorter one.}
\end{aligned}
\]
So the net cut leg $s_B$ is the derivative of $\Lambda_T - \Lambda_{T'}$ per unit of pooled mass removed: it is
the noise sensitivity of the premium mass the shorter clock removes from the pool, and it
carries the survivors' re-pricing. It is not the caught-only leg $\tilde s_B$. Under the extra clause of (C-2) it
splits into the two parts just named,
\begin{equation}
\label{eq:caught-split}
s_B \;=\; \tilde s_B \;-\; \frac{1-\varphi}{\varphi}\,\delta ,
\end{equation}
which part~(iii) proves.

The hypotheses are:
\begin{enumerate}[label=(C-\arabic*),leftmargin=3.2em,itemsep=3pt]
\item \emph{Fixed policies.} The plan menu, the execution policies and the cutoff vector are held
  fixed in $\kappa$ and held fixed across the two window margins, the threshold margin $\tau$ is
  common, the noise intensity $\kappa$ is common, and the timing carries no feedback from order
  flow into the blockholder's executed order path. Liquidity enters the primitives in one place
  only, the law of the ternary noise mark, which is standing condition (S10) of the partition
  subsection; fixed policies are standing condition (S11), cited separately here. Part~(ii)
  consumes both conditions.
\item \emph{Differentiability.} The two conditional expectations in \eqref{eq:caught-sens} are
  differentiable in $\kappa$ at the point compared. For the split \eqref{eq:caught-split} into
  $\tilde s_B$ and $\delta$, $\Eop[h^{(T)}\mid B]$ is also differentiable. The kernel is the
  pinned version of
  standing condition (S8), so $h^{(T)}$ and $h^{(T')}$ are well-defined random variables rather
  than equivalence classes; part~(iii) consumes that.
\item \emph{Non-degeneracy.} $\Prb(B) > 0$, $0 < \Omega(\tau,T)$, $\Omega(\tau,T') < 1$, and
  $s_A \ne 0$, equivalently $\mathcal{S}_P(\tau,T) > 0$. The clause $0 < \Omega(\tau,T)$ is there
  so that the clock theorem applies. Wherever this corollary is used to read the criterion
  $W_T C_T \le 1$ in Theorem~\ref{thm:clock}, that theorem's hypotheses also apply. They include
  the hypotheses of Proposition~\ref{prop:factorisation} at both clocks, including the invariance
  of the flagged endpoint in $\kappa$ at fixed policies.
\end{enumerate}

\begin{corollary}[Who gets caught]
\label{cor:caught}
Under (C-1) to (C-3):
\begin{enumerate}[label=(\roman*),leftmargin=2.4em,itemsep=3pt]
\item \emph{The cut is nested.} $\mathcal{C}_F(\tau,T)\subseteq\mathcal{C}_F(\tau,T')$, hence
  $B \subseteq A$ and $A\setminus B = \mathcal{C}_P(\tau,T')$, and the weight leg is
  \begin{equation}
  \label{eq:caught-WT}
  W_T \;=\; 1-\varphi \;\le\; 1 ,
  \qquad \varphi\in(0,1).
  \end{equation}
\item \emph{The cell masses are $\kappa$-free.}
  $\partial_\kappa\Prb(A) = \partial_\kappa\Prb(B) = 0$, so $\varphi$ does not depend on $\kappa$.
\item \emph{The cut identity.}
  \begin{equation}
  \label{eq:caught-identity}
  s_{A\setminus B}
  \;=\;
  \frac{\Prb(A)\,s_A - \Prb(B)\,s_B}{\Prb(A)-\Prb(B)}
  \;=\;
  \frac{s_A - \varphi\,s_B}{1-\varphi}
  \;=\;
  \frac{s_A - \varphi\,\tilde s_B}{1-\varphi} \;+\; \delta ,
  \end{equation}
  the last expression under the extra differentiability of (C-2). Equivalently
  $s_A = (1-\varphi)\,s_{A\setminus B} + \varphi\,s_B$: the longer clock's pooled sensitivity is
  the mass-weighted average of the sensitivity of what the shorter clock leaves in the pool and
  the net cut leg $s_B$, the derivative of $\Lambda_T-\Lambda_{T'}$ per unit of pooled mass removed, which
  carries the survivors' re-pricing.
\item \emph{The composition ratio.}
  \begin{equation}
  \label{eq:caught-CT}
  \begin{aligned}
  C_T \;&=\; \frac{|s_A-\varphi s_B|}{(1-\varphi)\,|s_A|}, \\[3pt]
  C_T \le 1 \quad&\Longleftrightarrow\quad
  (s_B-s_A)\bigl(\varphi s_B-(2-\varphi)s_A\bigr) \le 0 ,
  \end{aligned}
  \end{equation}
  that is, if and only if $s_B$ lies weakly between $s_A$ and
  $\bigl((2-\varphi)/\varphi\bigr)s_A$.
\item \emph{The attenuation criterion.}
  \begin{equation}
  \label{eq:caught-WTCT}
  \begin{aligned}
  W_T C_T \;&=\; \frac{|s_A-\varphi s_B|}{|s_A|}, \\[3pt]
  W_T C_T \le 1 \quad&\Longleftrightarrow\quad
  s_B\bigl(2s_A-\varphi s_B\bigr) \ge 0 ,
  \end{aligned}
  \end{equation}
  that is, if and only if $s_B$ lies weakly between $0$ and $(2/\varphi)s_A$.
\item \emph{Readings.} Write $\rho = s_B/s_A$, which is well defined by (C-3).
  \begin{enumerate}[label=(\alph*),leftmargin=2.2em,itemsep=2pt]
  \item If $s_A s_B \le 0$, including $s_B = 0$, then $C_T > 1$ strictly.
  \item If $s_A s_B > 0$, then
    $C_T \le 1$ if and only if
    $|s_A| \le |s_B| \le \bigl((2-\varphi)/\varphi\bigr)|s_A|$.
  \item If the cut does not reverse the pooled sensitivity, $s_{A\setminus B}\,s_A \ge 0$, then
    $C_T \le 1$ if and only if $\rho \ge 1$. The inequality $\rho \ge 1 > 0$ already forces a
    common sign, so that reading is $|s_B| \ge |s_A|$. Under a common sign the proviso
    $\sigma \ge 0$ is $|s_B| \le |s_A|/\varphi$.
  \item $W_T C_T \le 1$ if and only if $s_B$ shares the sign of $s_A$ or is zero, and
    $|s_B| \le (2/\varphi)|s_A|$.
  \item The two upper limits $(2-\varphi)/\varphi$ and $2/\varphi$ are strictly decreasing in
    $\varphi$ and diverge as $\varphi\downarrow0$. For every $M \ge 1$ and every cut with
    $\varphi \le 2/(1+M)$ and $|s_B| \le M|s_A|$, the upper limits in (b) and (d) are slack, so
    there $C_T \le 1$ reads: a common sign together with $|s_B| \ge |s_A|$; and $W_T C_T \le 1$
    reads: a common sign or $s_B = 0$.
  \end{enumerate}
\end{enumerate}
\end{corollary}

\begin{proof}
\emph{Part (i).} The disclosure indicator at window margin $T$ is
$D(\tau,T) = \ind\{a=1,\ c(\tau)+T\le H\}$, where $c(\tau)$ is the round in which the executed
stake path crosses $\tau$. If $D(\tau,T)=1$ then $c(\tau)+T' \le c(\tau)+T \le H$ because
$T'<T$, so $D(\tau,T')=1$. Hence $\mathcal{C}_F(\tau,T)\subseteq\mathcal{C}_F(\tau,T')$ and, taking
complements, $\mathcal{C}_P(\tau,T')\subseteq\mathcal{C}_P(\tau,T)=A$. Therefore
$B = \mathcal{C}_F(\tau,T')\cap A \subseteq A$ and
$A\setminus B = A\cap\mathcal{C}_P(\tau,T') = \mathcal{C}_P(\tau,T')$. Taking probabilities,
$\Omega(\tau,T) \le \Omega(\tau,T') < 1$, so both pooled cells carry positive mass; since
$B\subseteq A$, $\Prb(A\setminus B) = \Prb(A)-\Prb(B)$, and
\[
W_T \;=\; \frac{1-\Omega(\tau,T')}{1-\Omega(\tau,T)}
\;=\; \frac{\Prb(A\setminus B)}{\Prb(A)}
\;=\; 1-\varphi .
\]
By (C-3), $\Prb(B)>0$ gives $\varphi>0$ and $\Omega(\tau,T')<1$ gives $\Prb(A\setminus B)>0$, hence
$\varphi<1$.

\emph{Part (ii).} Under fixed policies the selected plan and the executed order path are functions
of the plan index and the signal alone, and by the no-feedback clause of (C-1) the order placed in
a round does not read realised order flow. So the crossing round $c(\tau)$, and with it
$D(\tau,T)$ and $D(\tau,T')$, are measurable with respect to the plan and signal. The intensity
$\kappa$ parametrises the noise law only; the laws of the value, the signal noise and the bidder
draw carry no $\kappa$, and (C-1) holds the cutoff vector fixed in $\kappa$, so the joint law of
the plan and the signal is one and the same at every $\kappa$. Hence $\Prb(A)$ and $\Prb(B)$ do not
vary with $\kappa$, and neither does their ratio $\varphi$. This is the same constancy the clock
theorem uses for the flagged weight, applied at both window margins.

\emph{Part (iii).} By part (i), $B$ and $A\setminus B$ partition $A$ and both have positive mass.
Write the pooled contributions to the expected premium at the two clocks, in kernel units, as
\[
\begin{aligned}
\Lambda_T \;&=\; \Eop\bigl[h^{(T)}\ind_A\bigr]
  \;=\; \Prb(A)\,\Eop\bigl[h^{(T)}\mid A\bigr], \\[3pt]
\Lambda_{T'} \;&=\; \Eop\bigl[h^{(T')}\ind_{A\setminus B}\bigr]
  \;=\; \Prb(A\setminus B)\,\Eop\bigl[h^{(T')}\mid A\setminus B\bigr] .
\end{aligned}
\]
By part (ii) the two masses are constants in $\kappa$, so the derivative that (C-2) supplies passes
through them and leaves them alone:
\[
\partial_\kappa \Lambda_T \;=\; \Prb(A)\,s_A ,
\qquad
\partial_\kappa \Lambda_{T'} \;=\; \bigl(\Prb(A)-\Prb(B)\bigr)\,s_{A\setminus B} .
\]
Definition \eqref{eq:caught-sB} reads
$s_B = \partial_\kappa(\Lambda_T - \Lambda_{T'})/\Prb(B)$, a division by a positive number, so
\[
\Prb(B)\,s_B \;=\; \Prb(A)\,s_A - \bigl(\Prb(A)-\Prb(B)\bigr)\,s_{A\setminus B} ,
\]
and solving for $s_{A\setminus B}$, which the positive mass $\Prb(A)-\Prb(B)$ permits, gives the
first equality of \eqref{eq:caught-identity}. Dividing numerator and denominator by $\Prb(A)$ gives
the second. Multiplying through by $1-\varphi$ and collecting terms gives
$s_A = (1-\varphi)s_{A\setminus B}+\varphi s_B$.

For the third expression, add and subtract $\Eop[h^{(T)}\ind_{A\setminus B}]$ inside
$\Lambda_T-\Lambda_{T'}$. Since $\ind_A = \ind_{A\setminus B}+\ind_B$ by part (i),
\[
\begin{aligned}
\Lambda_T - \Lambda_{T'}
\;&=\; \Eop\bigl[h^{(T)}\ind_B\bigr]
  \;-\; \Eop\bigl[\bigl(h^{(T')}-h^{(T)}\bigr)\ind_{A\setminus B}\bigr] \\[3pt]
\;&=\; \Prb(B)\,\Eop\bigl[h^{(T)}\mid B\bigr]
  \;-\; \Prb(A\setminus B)\,\Eop\bigl[h^{(T')}-h^{(T)}\mid A\setminus B\bigr] .
\end{aligned}
\]
Both masses are $\kappa$-free by part (ii), and the extra clause of (C-2) makes the first term
differentiable, hence so is the second. Differentiating and dividing by $\Prb(B)$ gives
$s_B = \tilde s_B - \bigl((1-\varphi)/\varphi\bigr)\delta$, which is \eqref{eq:caught-split}.
Substituting it into the second equality of \eqref{eq:caught-identity} gives
the third.

\emph{Part (iv).} By \eqref{eq:caught-SP} and part (i), $C_T = |s_{A\setminus B}|/|s_A|$, and by
\eqref{eq:caught-identity} with $1-\varphi>0$,
\[
C_T \;=\; \frac{|s_A-\varphi s_B|}{(1-\varphi)|s_A|} .
\]
The denominator is non-zero because $s_A \ne 0$ in (C-3). Both sides of $C_T\le 1$ are
non-negative, so $C_T\le 1$ if and only if $s_{A\setminus B}^2\le s_A^2$, and multiplying by
$(1-\varphi)^2>0$, if and only if $u^2-w^2\le 0$, where $u=s_A-\varphi s_B$ and
$w=(1-\varphi)s_A$. Since $u-w=\varphi(s_A-s_B)$ and $u+w=(2-\varphi)s_A-\varphi s_B$,
\[
u^2-w^2 \;=\; (u-w)(u+w) \;=\; \varphi\,(s_A-s_B)\bigl((2-\varphi)s_A-\varphi s_B\bigr).
\]
Since $\varphi>0$, the sign of the whole is the sign of
$(s_A-s_B)\bigl((2-\varphi)s_A-\varphi s_B\bigr)$, and negating both factors leaves the product
unchanged, which gives the stated form
$(s_B-s_A)\bigl(\varphi s_B-(2-\varphi)s_A\bigr)\le 0$. Writing that product as
$\varphi(s_B-s_A)\bigl(s_B-\bigl((2-\varphi)/\varphi\bigr)s_A\bigr)$ shows it is non-positive
exactly when $s_B$ lies weakly between the two roots $s_A$ and
$\bigl((2-\varphi)/\varphi\bigr)s_A$.

\emph{Part (v).} By part (i) and part (iv),
\[
W_T C_T \;=\; (1-\varphi)\,\frac{|s_{A\setminus B}|}{|s_A|}
\;=\; \frac{\bigl|(1-\varphi)s_{A\setminus B}\bigr|}{|s_A|}
\;=\; \frac{|s_A-\varphi s_B|}{|s_A|} .
\]
As in part (iv), $W_T C_T\le 1$ if and only if $u^2-s_A^2\le 0$ for $u=s_A-\varphi s_B$, and since
$u-s_A=-\varphi s_B$ and $u+s_A=2s_A-\varphi s_B$,
\[
u^2-s_A^2 \;=\; (u-s_A)(u+s_A) \;=\; -\varphi\,s_B\bigl(2s_A-\varphi s_B\bigr),
\]
so the criterion is $s_B(2s_A-\varphi s_B)\ge 0$. Writing it as
$\varphi\,s_B\bigl((2/\varphi)s_A-s_B\bigr)\ge 0$ shows it holds exactly when $s_B$ lies weakly
between $0$ and $(2/\varphi)s_A$.

\emph{Part (vi).} Since $\varphi\in(0,1)$, the factor $(2-\varphi)/\varphi$ is strictly positive,
so the two roots in part (iv) are non-zero and share the sign of $s_A$, and the closed interval
they bound does not contain zero. If $s_As_B\le 0$ then $s_B$ is zero or lies on the far side of
zero from both roots, so it is strictly outside that interval, the product in
\eqref{eq:caught-CT} is strictly positive, and $C_T>1$. That is (a). For (b), divide the interval
condition by $s_A$: with a common sign, $\rho>0$ and $s_B$ between $s_A$ and
$\bigl((2-\varphi)/\varphi\bigr)s_A$ reads $1\le\rho\le(2-\varphi)/\varphi$, and multiplying by
$|s_A|$ gives the stated magnitudes. For (c), set
$\sigma = s_{A\setminus B}/s_A = (1-\varphi\rho)/(1-\varphi)$ by \eqref{eq:caught-identity}. The
proviso $s_{A\setminus B}s_A\ge0$ is $\sigma\ge0$, so $C_T = |\sigma| = \sigma$, and $C_T\le1$ if
and only if $1-\varphi\rho\le1-\varphi$, that is $\rho\ge1$. The inequality $\rho\ge1>0$ already
forces a common sign, under which $\rho\ge1$ is $|s_B|\ge|s_A|$. Under a common sign the proviso
$\sigma\ge0$ is $\varphi\rho\le1$, that is $|s_B|\le|s_A|/\varphi$. For (d), divide the interval
condition of part (v) by $s_A$: it reads $0\le\rho\le2/\varphi$, which is a common sign or
$s_B=0$, together with $|s_B|\le(2/\varphi)|s_A|$. For (e), both $(2-\varphi)/\varphi$ and
$2/\varphi$ are strictly decreasing on $(0,1)$ and tend to $+\infty$ as $\varphi\downarrow0$. Fix
$M\ge1$ and $\varphi\le2/(1+M)$. Then $(2-\varphi)/\varphi\ge M$, because that inequality is
$2-\varphi\ge M\varphi$, and $2/\varphi\ge(2-\varphi)/\varphi\ge M$. So $|s_B|\le M|s_A|$ makes
the upper limits in (b) and (d) slack, leaving in (b) the common sign with $|s_B|\ge|s_A|$ and in
(d) the common sign or $s_B=0$.
\end{proof}

\paragraph{Reading.}
Part~(iii) is the whole content in one line: the longer clock's pooled sensitivity is an average
over what the shorter clock leaves in the pool and what it takes out of the pool, with weights the
two masses, and those weights do not move with liquidity. So shortening the clock pulls the pooled
sensitivity down only if what it takes out sat above the average. Parts~(iv) and~(vi) put the
bound on how far above. What the clock removes must move with the pool in $\kappa$ and be at least
as noise-sensitive as the pool, and it must not be so much more sensitive that the remaining pool
is driven through zero to a larger magnitude than the pool started with. A cut that takes a small
share of the pool faces only the first two requirements when, as part~(vi)(e) states, there is
some $M\ge 1$ with $\varphi \le 2/(1+M)$ and $|s_B|\le M|s_A|$. Part~(v) does the same for the criterion
the clock theorem states, $W_T C_T\le1$, where the weight leg is already working in the same
direction: there what the clock removes need only move with the pool and stay inside $2/\varphi$
times its sensitivity.

What the clock removes has two parts, and \eqref{eq:caught-split} separates them. The first is the
caught-only leg, the noise sensitivity of the newly caught histories read at the prices the longer
clock gave them. The second is the re-pricing term: the surviving pool is a smaller pool once the
clock shortens, and the market reads it against that smaller pool, so its own sensitivity moves
too. The criterion runs on the sum, which is the sensitivity of the premium mass that leaves the
pool. When the re-pricing term is zero the sum is the first part alone, and the criterion is
a statement about the caught histories and nothing else.

The corollary turns the ratio $C_T$ into a question with an answer at every calibration node, and
that question is who gets caught. Under a common sign, $s_A s_B > 0$, the composition ratio
satisfies $C_T \le 1$ if and only if the net cut leg $s_B$ lies weakly between $s_A$ and
$\bigl((2-\varphi)/\varphi\bigr)\,s_A$, which is part~(iv); and the shorter clock lowers noise
sensitivity overall, $W_T C_T \le 1$, if and only if $s_B$ lies weakly between $0$ and
$(2/\varphi)\,s_A$, which is part~(v). The two bands differ: the weight leg lets the overall
criterion accept a net cut leg less sensitive than the pool, which the composition criterion alone
does not.

% =============================================================================
% The same accounting law at the threshold margin.
% =============================================================================

\csname unlabelledtrue\endcsname
\begin{corollary}[Who gets caught at the threshold margin]
\label{cor:caught-threshold}
Fix a common window margin $T$ and liquidity $\kappa$, and compare $b_0<\tau'<\tau$, with
$(\tau',T)$ the tighter rule. Write
\begin{equation}
\label{eq:caught-tau-sets}
A_\tau=\mathcal C_P(\tau,T),
\qquad
B_\tau=\mathcal C_F(\tau',T)\setminus\mathcal C_F(\tau,T),
\qquad
R_\tau=A_\tau\setminus B_\tau,
\end{equation}
and let $h^{(\tau)}$ and $h^{(\tau')}$ be the kernels supplied by the two rules. Set
\begin{align}
\label{eq:caught-tau-definitions}
\varphi_\tau&=\frac{\Prb(B_\tau)}{\Prb(A_\tau)},
&s_{A,\tau}&=\partial_\kappa\Eop[h^{(\tau)}\mid A_\tau],
&s_{R,\tau}&=\partial_\kappa\Eop[h^{(\tau')}\mid R_\tau],
\nonumber\\
\widetilde s_{B,\tau}&=\partial_\kappa\Eop[h^{(\tau)}\mid B_\tau],
&\delta_\tau&=\partial_\kappa\Eop[h^{(\tau')}-h^{(\tau)}\mid R_\tau],
\end{align}
and define the net cut leg by
\begin{equation}
\label{eq:caught-tau-sB}
s_{B,\tau}:=\frac{1}{\Prb(B_\tau)}\partial_\kappa\left(
\Eop[h^{(\tau)}\ind_{A_\tau}]
-\Eop[h^{(\tau')}\ind_{R_\tau}]
\right).
\end{equation}
Assume standing conditions (S1) to (S11) and the following conditions.
\begin{enumerate}[label=(C$\tau$-\arabic*),leftmargin=3.8em,itemsep=3pt]
\item \emph{A fixed threshold cut.} The window margin and liquidity are common, and
  $b_0<\tau'<\tau$. Conditions (S4), (S5), and (S7) apply at both thresholds, so both disclosure
  indicators are functions of the same plan and signal. Liquidity enters only through the noise law
  by (S10), and the policies are common by (S11).
\item \emph{Differentiability and pinned kernels.} The maps
  $\kappa\mapsto\Eop[h^{(\tau)}\mid A_\tau]$ and
  $\kappa\mapsto\Eop[h^{(\tau')}\mid R_\tau]$ are differentiable at the compared liquidity.
  Condition (S8) pins both kernels. For the split into the caught-only leg and the re-pricing term,
  $\Eop[h^{(\tau)}\mid B_\tau]$ is also differentiable.
\item \emph{Non-degeneracy and factorisation.}
  $0<\Prb(B_\tau)<\Prb(A_\tau)$, $0<\Omega(\tau,T)$,
  $\Omega(\tau',T)<1$, and $s_{A,\tau}\ne0$. Whenever the aggregate criterion is invoked, the
  hypotheses of Proposition~\ref{prop:factorisation}, including flagged-endpoint invariance, hold
  at both thresholds.
\item \emph{Erasure representation.} Only the identification with
  Condition~\ref{cond:D} requires standing conditions (S12) and (S13), order size two, and
  $\kappa\in(0,1)$, so that Lemmas~\ref{lem:g1} to~\ref{lem:g3} apply. The cut algebra below does
  not use this clause.
\end{enumerate}
Under (C$\tau$-1) to (C$\tau$-3):
\begin{enumerate}[label=(\roman*),leftmargin=2.4em,itemsep=3pt]
\item \emph{The cut is nested.}
  $\mathcal C_F(\tau,T)\subseteq\mathcal C_F(\tau',T)$, hence
  $B_\tau\subseteq A_\tau$ and $R_\tau=\mathcal C_P(\tau',T)$, and
  \begin{equation}
  \label{eq:caught-tau-W}
  W_\tau=\frac{1-\Omega(\tau',T)}{1-\Omega(\tau,T)}
  =1-\varphi_\tau\le1,
  \qquad \varphi_\tau\in(0,1).
  \end{equation}
\item \emph{The cell masses are $\kappa$-free.}
  $\partial_\kappa\Prb(A_\tau)=\partial_\kappa\Prb(B_\tau)=0$, so
  $\varphi_\tau$ does not depend on $\kappa$.
\item \emph{The cut identity.}
  \begin{equation}
  \label{eq:caught-tau-identity}
  s_{R,\tau}
  =\frac{\Prb(A_\tau)s_{A,\tau}-\Prb(B_\tau)s_{B,\tau}}
  {\Prb(A_\tau)-\Prb(B_\tau)}
  =\frac{s_{A,\tau}-\varphi_\tau s_{B,\tau}}{1-\varphi_\tau}
  =\frac{s_{A,\tau}-\varphi_\tau\widetilde s_{B,\tau}}{1-\varphi_\tau}
  +\delta_\tau,
  \end{equation}
  where the last expression uses the extra differentiability in (C$\tau$-2). Equivalently,
  \begin{equation}
  \label{eq:caught-tau-average}
  s_{A,\tau}=(1-\varphi_\tau)s_{R,\tau}+\varphi_\tau s_{B,\tau},
  \qquad
  s_{B,\tau}=\widetilde s_{B,\tau}
  -\frac{1-\varphi_\tau}{\varphi_\tau}\delta_\tau.
  \end{equation}
\item \emph{The composition ratio.}
  \begin{equation}
  \label{eq:caught-tau-C}
  \begin{aligned}
  C_\tau&=\frac{|s_{A,\tau}-\varphi_\tau s_{B,\tau}|}
  {(1-\varphi_\tau)|s_{A,\tau}|},\\
  C_\tau\le1
  &\quad\Longleftrightarrow\quad
  (s_{B,\tau}-s_{A,\tau})
  \bigl(\varphi_\tau s_{B,\tau}-(2-\varphi_\tau)s_{A,\tau}\bigr)\le0.
  \end{aligned}
  \end{equation}
  Thus $C_\tau\le1$ exactly when the net cut leg lies weakly between $s_{A,\tau}$ and
  $((2-\varphi_\tau)/\varphi_\tau)s_{A,\tau}$. Under (C$\tau$-4), this is equivalent to
  Condition~\ref{cond:D}.
\item \emph{The attenuation criterion.}
  \begin{equation}
  \label{eq:caught-tau-WC}
  \begin{aligned}
  W_\tau C_\tau
  &=\frac{|s_{A,\tau}-\varphi_\tau s_{B,\tau}|}{|s_{A,\tau}|},\\
  W_\tau C_\tau\le1
  &\quad\Longleftrightarrow\quad
  s_{B,\tau}(2s_{A,\tau}-\varphi_\tau s_{B,\tau})\ge0.
  \end{aligned}
  \end{equation}
  Thus attenuation holds exactly when the net cut leg lies weakly between $0$ and
  $(2/\varphi_\tau)s_{A,\tau}$. Under (C$\tau$-3), this is equivalent to the tighter threshold
  weakly lowering aggregate liquidity sensitivity.
\item \emph{Readings.} Write $\rho_\tau=s_{B,\tau}/s_{A,\tau}$, which is well defined by
  (C$\tau$-3).
  \begin{enumerate}[label=(\alph*),leftmargin=2.2em,itemsep=2pt]
  \item If $s_{A,\tau}s_{B,\tau}\le0$, including $s_{B,\tau}=0$, then $C_\tau>1$.
  \item If $s_{A,\tau}s_{B,\tau}>0$, then $C_\tau\le1$ if and only if
    $|s_{A,\tau}|\le|s_{B,\tau}|\le
    ((2-\varphi_\tau)/\varphi_\tau)|s_{A,\tau}|$.
  \item If $s_{R,\tau}s_{A,\tau}\ge0$, then $C_\tau\le1$ if and only if
    $\rho_\tau\ge1$. That inequality forces $s_{B,\tau}$ and $s_{A,\tau}$ to have a common sign and
    says $|s_{B,\tau}|\ge|s_{A,\tau}|$. Under that common sign, the proviso
    $s_{R,\tau}s_{A,\tau}\ge0$ is
    $|s_{B,\tau}|\le|s_{A,\tau}|/\varphi_\tau$.
  \item $W_\tau C_\tau\le1$ if and only if $s_{B,\tau}$ has the sign of
    $s_{A,\tau}$ or is zero, and
    $|s_{B,\tau}|\le(2/\varphi_\tau)|s_{A,\tau}|$.
  \item The upper limits $(2-\varphi_\tau)/\varphi_\tau$ and $2/\varphi_\tau$ fall as
    $\varphi_\tau$ rises and diverge as $\varphi_\tau\downarrow0$. If $M\ge1$,
    $\varphi_\tau\le2/(1+M)$, and $|s_{B,\tau}|\le M|s_{A,\tau}|$, the upper limits are slack.
    The composition condition then asks for a common sign and
    $|s_{B,\tau}|\ge|s_{A,\tau}|$. The aggregate condition asks only for a common sign or
    $s_{B,\tau}=0$.
  \end{enumerate}
\end{enumerate}
\end{corollary}
\csname unlabelledfalse\endcsname

\begin{proof}
By (C$\tau$-1) and Lemma~\ref{lem:threshold-weight}, at either threshold the disclosure indicator
has the product form
\[
D(\tau_0,T)=\ind\{a=1\}\ind\{B(s,H-T)\ge\tau_0\},
\qquad \tau_0\in\{\tau',\tau\}.
\]
If a history is flagged at $\tau$, then its common stake path satisfies
$B(s,H-T)\ge\tau>\tau'$, so it is flagged at $\tau'$. This proves the flagged-set inclusion.
Taking complements gives $R_\tau=\mathcal C_P(\tau',T)$. The mass identity then gives
\eqref{eq:caught-tau-W}, and (C$\tau$-3) gives $0<\varphi_\tau<1$.

Conditions (S3), (S7), (S10), and (S11) make $A_\tau$ and $B_\tau$ events determined by the common plan
and signal, whose law does not move with $\kappa$. Their masses and their ratio are therefore
$\kappa$-free. This proves part~(ii).

For part~(iii), put
\[
\Lambda_\tau=\Eop[h^{(\tau)}\ind_{A_\tau}],
\qquad
\Lambda_{\tau'}=\Eop[h^{(\tau')}\ind_{R_\tau}].
\]
Part~(ii) and (C$\tau$-2) give
\[
\partial_\kappa\Lambda_\tau=\Prb(A_\tau)s_{A,\tau},
\qquad
\partial_\kappa\Lambda_{\tau'}=\Prb(R_\tau)s_{R,\tau}.
\]
Definition~\eqref{eq:caught-tau-sB} and
$\Prb(R_\tau)=\Prb(A_\tau)-\Prb(B_\tau)$ now give the first two equalities in
\eqref{eq:caught-tau-identity} and the average identity in
\eqref{eq:caught-tau-average}. For the split, use
$\ind_{A_\tau}=\ind_{B_\tau}+\ind_{R_\tau}$ to write
\[
\Lambda_\tau-\Lambda_{\tau'}
=\Prb(B_\tau)\Eop[h^{(\tau)}\mid B_\tau]
-\Prb(R_\tau)\Eop[h^{(\tau')}-h^{(\tau)}\mid R_\tau].
\]
Differentiate and divide by $\Prb(B_\tau)$. This proves the second identity in
\eqref{eq:caught-tau-average}; substituting it into the average identity proves the last equality in
\eqref{eq:caught-tau-identity}.

By (S9) and Lemma~\ref{lem:two-cell},
\[
\mathcal S_P(\kappa,\tau,T)=\Delta_m|s_{A,\tau}|,
\qquad
\mathcal S_P(\kappa,\tau',T)=\Delta_m|s_{R,\tau}|.
\]
Use \eqref{eq:caught-tau-identity} and $1-\varphi_\tau>0$ to obtain the first line of
\eqref{eq:caught-tau-C}. To test $C_\tau\le1$, set
$u=s_{A,\tau}-\varphi_\tau s_{B,\tau}$ and
$w=(1-\varphi_\tau)s_{A,\tau}$. Then
\[
u^2-w^2
=\varphi_\tau(s_{A,\tau}-s_{B,\tau})
\bigl((2-\varphi_\tau)s_{A,\tau}-\varphi_\tau s_{B,\tau}\bigr).
\]
The denominator in $C_\tau$ is positive. Squaring therefore preserves the comparison, and the
display is nonpositive exactly when $s_{B,\tau}$ lies between the two roots in part~(iv). Under
(C$\tau$-4), Lemma~\ref{lem:g3} identifies the two pooled sensitivities with the absolute revelation
values, so this is Condition~\ref{cond:D}.

Multiplying the first line of \eqref{eq:caught-tau-C} by
$W_\tau=1-\varphi_\tau$ proves the first line of \eqref{eq:caught-tau-WC}. The second follows from
\[
(s_{A,\tau}-\varphi_\tau s_{B,\tau})^2-s_{A,\tau}^2
=-\varphi_\tau s_{B,\tau}(2s_{A,\tau}-\varphi_\tau s_{B,\tau}).
\]
The two roots are $0$ and $(2/\varphi_\tau)s_{A,\tau}$. The factorisation assumed in
(C$\tau$-3) turns $W_\tau C_\tau$ into the ratio of aggregate sensitivities, proving part~(v).

For part~(vi), the two roots in part~(iv) are nonzero and have the sign of $s_{A,\tau}$, so zero
lies strictly outside the closed interval between them. This proves (a). Under a common sign,
division by $s_{A,\tau}$ turns their interval into
$1\le\rho_\tau\le(2-\varphi_\tau)/\varphi_\tau$, proving (b). The average identity gives
\[
\frac{s_{R,\tau}}{s_{A,\tau}}
=\frac{1-\varphi_\tau\rho_\tau}{1-\varphi_\tau}.
\]
Under the proviso in (c), the left side is nonnegative, so its absolute value is at most one exactly
when $\rho_\tau\ge1$; the same display gives the stated upper limit. Dividing the interval in
part~(v) by $s_{A,\tau}$ proves (d). Finally, both upper limits fall in $\varphi_\tau$ and diverge at
zero. The inequality $\varphi_\tau\le2/(1+M)$ implies
$(2-\varphi_\tau)/\varphi_\tau\ge M$ and $2/\varphi_\tau\ge M$, which proves (e).
\end{proof}

% =============================================================================
% The split and the polynomial implications for either dial.
% =============================================================================

\csname unlabelledtrue\endcsname
\begin{lemma}[The cut split at both margins]
\label{lem:cut-split}
Assume standing conditions (S1) to (S11). Compare either a shorter clock $T'<T$ at common
$(\tau,\kappa)$ or a tighter threshold $b_0<\tau'<\tau$ at common $(T,\kappa)$. Call the looser rule
$r_-$ and the tighter rule $r_+$, and write
\[
A=\mathcal C_{P,r_-},
\qquad
B=\mathcal C_{F,r_+}\setminus\mathcal C_{F,r_-},
\qquad
R=A\setminus B,
\qquad
\varphi=\frac{\Prb(B)}{\Prb(A)}.
\]
Here $R$ is the surviving tighter pool; Lemma~\ref{lem:silence} uses $A$ for that role.
Let $h^-$ and $h^+$ be the kernels supplied by $r_-$ and $r_+$. The following conditions are in
force.
\begin{enumerate}[label=(K-\arabic*),leftmargin=3.2em,itemsep=3pt]
\item \emph{A non-null nested cut.} Conditions (S4), (S5), (S7), and (S11), with
  Corollary~\ref{cor:caught}(i) for the clock or Corollary~\ref{cor:caught-threshold}(i) for the
  threshold, give $R=\mathcal C_{P,r_+}$ and the displayed nesting at the chosen margin, and
  $0<\Prb(B)<\Prb(A)$.
\item \emph{Pinned differentiable terms.} Condition (S8) pins $h^-$ and $h^+$. The maps
  $\kappa\mapsto\Eop[h^-\mid B]$ and
  $\kappa\mapsto\Eop[h^- -h^+\mid R]$ are differentiable. Conditions (S1) to (S3), (S7),
  (S10), and (S11) make the masses of $A$, $B$, and $R$ independent of $\kappa$.
\end{enumerate}
Define the caught-only leg, the re-pricing term oriented with the net cut, and the net cut leg by
\begin{align}
\label{eq:cut-split-definitions}
\widetilde s_B&=\partial_\kappa\Eop[h^-\mid B],
&\delta_{\mathrm{cut}}&=\partial_\kappa\Eop[h^- -h^+\mid R],
\nonumber\\
s_B&=\frac{1}{\Prb(B)}\partial_\kappa
\left(\Eop[h^-\ind_A]-\Eop[h^+\ind_R]\right).
\end{align}
Thus $\delta_{\mathrm{cut}}=-\delta$ for the clock notation above and
$\delta_{\mathrm{cut}}=-\delta_\tau$ for the threshold notation above. This orientation makes both
contributions enter the net cut with positive coefficients.
Then, at either margin,
\begin{equation}
\label{eq:cut-split-both}
s_B=\widetilde s_B+\frac{1-\varphi}{\varphi}\,\delta_{\mathrm{cut}}.
\end{equation}
The caught-only leg prices the newly flagged histories under the looser rule. The re-pricing term
prices the survivors under both rules.
\end{lemma}
\csname unlabelledfalse\endcsname

\begin{proof}
The sets $B$ and $R$ partition $A$, so
\[
\begin{aligned}
\Eop[h^-\ind_A]-\Eop[h^+\ind_R]
&=\Eop[h^-\ind_B]+\Eop[(h^- -h^+)\ind_R]\\
&=\Prb(B)\Eop[h^-\mid B]
  +\Prb(R)\Eop[h^- -h^+\mid R].
\end{aligned}
\]
The three masses are $\kappa$-free under (K-2). Differentiate, divide by $\Prb(B)>0$, and use
$\Prb(R)/\Prb(B)=(1-\varphi)/\varphi$. This gives
\eqref{eq:cut-split-both}.
\end{proof}

\csname unlabelledtrue\endcsname
\begin{lemma}[One crossing decides the composition band]
\label{lem:one-crossing}
Take a non-null nested clock or threshold cut at fixed policies and an interior liquidity
$\kappa\in(0,1)$ in the order-size-two model. Let $r_-$ be the looser rule and
$r_+$ the tighter rule, and write
\[
A=\mathcal C_{P,r_-},
\qquad
B=\mathcal C_{F,r_+}\setminus\mathcal C_{F,r_-},
\qquad
R=A\setminus B=\mathcal C_{P,r_+},
\qquad
\varphi=\frac{\Prb(B)}{\Prb(A)}\in(0,1).
\]
Assume standing conditions (S1) to (S13), so
Lemma~\ref{lem:g3} represents each pooled premium as a polynomial in $\kappa$ with coefficients that
do not depend on $\kappa$. In kernel units write
\begin{equation}
\label{eq:one-crossing-slopes}
s^q(\kappa)=-\frac12\sum_{k=0}^{H}(1-\varepsilon)^k\varepsilon^{H-k}c_k^q,
\qquad \varepsilon=\kappa/2,
\qquad q\in\{-,+\},
\end{equation}
where
$s^-=\partial_\kappa\Eop[h^-\mid A]$ and
$s^+=\partial_\kappa\Eop[h^+\mid R]$ are the signed pooled slopes under $r_-$ and $r_+$. Set
$d_k=c_k^+-c_k^-$ and, for $v=(v_0,\ldots,v_H)$, define the reversed polynomial
\begin{equation}
\label{eq:one-crossing-polynomial}
P_v(x)=\sum_{k=0}^{H}v_kx^{H-k},
\qquad x=x(\kappa)=\frac{\kappa}{2-\kappa}.
\end{equation}
Assume:
\begin{enumerate}[label=(R\arabic*),leftmargin=3.0em,itemsep=3pt]
\item \emph{One crossing for each pooled cell.} After zero coefficients are deleted, each of
  $(c_0^-,\ldots,c_H^-)$ and $(c_0^+,\ldots,c_H^+)$ changes sign exactly once, from negative to
  positive. The actual endpoint coefficients satisfy $c_0^-<0<c_H^-$ and $c_0^+<0<c_H^+$.
\item \emph{One crossing for the survivor difference.} After zero coefficients are deleted,
  $(d_0,\ldots,d_H)$ changes sign exactly once, from positive to negative, and
  $d_0>0>d_H$.
\item \emph{Liquidity is above the three revelation roots.} Let $r_A$, $r_R$, and $r_D$ be the
  positive roots of $P_{c^-}$, $P_{c^+}$, and $P_d$, which (R1), (R2), and Descartes' rule make
  unique. At the liquidity compared,
  \[
  x(\kappa)>r_*:=\max\{r_A,r_R,r_D\}.
  \]
\end{enumerate}
Then
\begin{equation}
\label{eq:one-crossing-order}
0<s^+(\kappa)<s^-(\kappa),
\end{equation}
and the net cut leg defined by
$s^-=(1-\varphi)s^++\varphi s_B$ satisfies
\begin{equation}
\label{eq:one-crossing-band}
s^-<s_B<\frac{s^-}{\varphi}
<\frac{2-\varphi}{\varphi}s^-.
\end{equation}
Thus the net cut leg lies in the composition band and the tighter pool has strictly lower liquidity
sensitivity. The composition ratio at the chosen margin, as defined in
Corollary~\ref{cor:caught} or Corollary~\ref{cor:caught-threshold}, is
$C_{\mathrm{cut}}=|s^+|/|s^-|<1$.
\end{lemma}
\csname unlabelledfalse\endcsname

\begin{proof}
The coefficient list of each $P_{c^q}$ has exactly one sign change. Descartes' rule of signs gives
exactly one positive root, counted with multiplicity. The strict endpoint coefficients give
$P_{c^q}(0)=c_H^q>0$ and a negative leading coefficient, so
$P_{c^q}(x)<0$ above its positive root. The same argument applied to $P_d$ gives one positive root,
with $P_d(0)=d_H<0$ and a positive leading coefficient, so $P_d(x)>0$ above that root.

\begin{equation}
\label{eq:one-crossing-factor}
\sum_{k=0}^{H}(1-\varepsilon)^k\varepsilon^{H-k}v_k
=(1-\varepsilon)^H P_v(x(\kappa)).
\end{equation}
The factor on the right is positive. Since $P_d=P_{c^+}-P_{c^-}$ by linearity of $v\mapsto P_v$,
condition (R3) therefore gives
\[
\Wrev^-<\Wrev^+<0,
\qquad
\Wrev^q:=\sum_{k=0}^{H}(1-\varepsilon)^k\varepsilon^{H-k}c_k^q.
\]
Equation~\eqref{eq:one-crossing-slopes} proves
\eqref{eq:one-crossing-order}. Solving the cut identity for the net cut leg gives
\[
s_B=\frac{s^--(1-\varphi)s^+}{\varphi}.
\]
Hence
\[
s_B-s^-=\frac{1-\varphi}{\varphi}(s^--s^+)>0,
\qquad
\frac{s^-}{\varphi}-s_B=\frac{1-\varphi}{\varphi}s^+>0.
\]
Since $0<\varphi<1$ and $s^->0$,
$s^-/\varphi<((2-\varphi)/\varphi)s^-$. This proves the band and the composition conclusion.
\end{proof}

\csname unlabelledtrue\endcsname
\begin{lemma}[Distinct crossings reverse total sensitivity]
\label{lem:reversal}
Compare a non-null nested pair at fixed policies, with $r_-$ the looser rule and $r_+$ the tighter
rule. Let $s^-$ be the signed pooled slope under $r_-$ and let $s^+$ be the slope under $r_+$,
using its own kernel on $\mathcal C_{P,r_+}$. Write
$\Omega_q=\Prb(\mathcal C_{F,r_q})$ for $q\in\{-,+\}$. Assume:
\begin{enumerate}[label=(V-\arabic*),leftmargin=3.2em,itemsep=3pt]
\item \emph{Standing conditions and factorisation.} Conditions (S1) to (S13) hold,
  $0<\Omega_-<1$ and $0<\Omega_+<1$, and the hypotheses of
  Proposition~\ref{prop:factorisation} hold under both rules on an open set in $(0,1)$ containing
  every zero of the two signed pooled slopes.
\item \emph{Distinct one-crossing roots.} The functions $s^-$ and $s^+$ are continuous on $(0,1)$.
  Each has exactly one zero there and changes sign at it. Write the roots as $\kappa_-$ and
  $\kappa_+$ respectively, and assume $\kappa_-\ne\kappa_+$.
\end{enumerate}
Then there is an open interval $I$ containing the looser pool's root $\kappa_-$ on which the tighter
rule is strictly more sensitive in total:
\begin{equation}
\label{eq:reversal}
\mathcal S(\kappa,r_+)>\mathcal S(\kappa,r_-)
\qquad\text{for every }\kappa\in I.
\end{equation}
\end{lemma}
\csname unlabelledfalse\endcsname

\begin{proof}
Factorisation, (S9), and Step~3 of Proposition~\ref{prop:factorisation}, which makes $\Omega_q$
constant in $\kappa$, give, for $q\in\{-,+\}$,
\[
\mathcal S(\kappa,r_q)=(1-\Omega_q)\Delta_m|s^q(\kappa)|.
\]
At the looser pool's root, $s^-(\kappa_-)=0$. The roots are distinct and $s^+$ has no other zero,
so $s^+(\kappa_-)\ne0$. Since $1-\Omega_+>0$ and $\Delta_m>0$,
\[
\mathcal S(\kappa_-,r_+)>0=\mathcal S(\kappa_-,r_-).
\]
Both sides are continuous in $\kappa$. Their strict difference at $\kappa_-$ therefore remains
positive on some open interval containing $\kappa_-$ and contained in the neighbourhood from
(V-1). This proves \eqref{eq:reversal}.
\end{proof}


% proofs/09_regret.tex
% The benchmark-policy one-step deviation bound.
% Fragment, to be \input by appendix.tex.

\section{The benchmark policy: the interim regret bound}
\label{sec:regret}

At each calibration node, a one-step deviation changes the blockholder's plan at one signal while
the benchmark price system and beliefs stay fixed. The result below bounds the largest gain from
that deviation on the stated signal support.

\csname unlabelledtrue\endcsname
\begin{proposition}[Benchmark-policy interim regret]
\label{prop:regret}
Assume standing conditions (S1) to (S11). The fixed pass in (R-3) also satisfies the common noise
channel, control-node convention, and identified flagged types in (S12) to (S14). The following
additional conditions are in force.
\begin{enumerate}[label=(R-\arabic*),leftmargin=3.2em,itemsep=3pt]
\item \emph{Calibration.} The benchmark cutoffs are
  \[
  k^b=(0.9425017266871091, 1.8484512098302512),
  \]
  liquidity is $\kappa=0.5$, order size is two noise lumps, and $H=10$. The nodes pair
  $T\in\{3,5,10\}$ with
  \[
  \tau\in\left\{
  \begin{aligned}
  &0.09239820387429526, 0.09346755804663053, 0.09453534811956685,\\
  &0.09565657882778708, 0.09703177146201895
  \end{aligned}
  \right\},
  \]
  and the remaining primitives take their benchmark calibration values. In particular,
  $m_1>0$, $\bar b>b_0$, $\sigma_s>0$, $\sigma_\xi>0$, $\beta>0$, and $\chi>0$.
\item \emph{Signal support.} The signal has its calibrated Gaussian law restricted to
  $I_s=[s_{\mathrm{lo}},s_{\mathrm{hi}}]$, where
  $s_{\mathrm{lo}}=\mu_v-6\sigma_s$ and $s_{\mathrm{hi}}=\mu_v+6\sigma_s$. Its density is divided by the
  truncation constant
  \[
  Z_s=\Prb(s\in I_s)
  =\Phi\!\left(\frac{s_{\mathrm{hi}}-\mu_v}{\sigma_s}\right)
  -\Phi\!\left(\frac{s_{\mathrm{lo}}-\mu_v}{\sigma_s}\right)>0.
  \]
  The density is continuous, so every finite breakpoint set is null.
\item \emph{Fixed pooled pass and deviation convention.} At each node, the pooled pricing pass is
  computed once at the benchmark cutoffs with the run-up path present and is then held fixed. A
  deviation changes only the deviating signal's plan. Benchmark prices and beliefs do not react.
  Every off-path pooled history uses the fixed reference belief of the mark-path type that can
  produce it. At a cutoff the benchmark assigns Hold at the Exit-Hold cutoff and Voice at the
  Hold-Voice cutoff, so Exit is assigned for $s<k_1^b$, Hold for
  $k_1^b\le s<k_2^b$, and Voice for $s\ge k_2^b$.
\end{enumerate}

For calibration node $n$, let $U_{j,n}(s)$ be the payoff from plan
$j\in\{\mathsf{E},\mathsf{H},\mathsf{V}\}$ under (R-3), and let $j^b(s)$ be the benchmark plan. Define
\begin{equation}
\label{eq:regret-definition}
R_n(s)=\max_{j\in\{\mathsf{E},\mathsf{H},\mathsf{V}\}}U_{j,n}(s)-U_{j^b(s),n}(s).
\end{equation}
Use the benchmark payoff at the Hold-Voice cutoff as the payoff normaliser,
\begin{equation}
\label{eq:regret-normaliser}
N_n:=U_{\mathsf{V},n}(k_2^b)>0.
\end{equation}
The reported certificate value $\overline R_n$ satisfies
\begin{equation}
\label{eq:regret-bound}
\operatorname*{ess\,sup}_{s\in I_s}R_n(s)\le\overline R_n,
\qquad
\operatorname*{ess\,sup}_{s\in I_s}\frac{R_n(s)}{N_n}
\le\frac{\overline R_n}{N_n}.
\end{equation}
The record reports the bound at every node.
\end{proposition}
\csname unlabelledfalse\endcsname

\begin{proof}
Fix a node and suppress its subscript. Form a partition of $I_s$ from the two benchmark cutoffs,
every jump of the Voice accumulation length $n(s)$ in $I_s$, every signal in $I_s$ at which a
threshold-crossing date changes, and the two support endpoints. Remove exact duplicates and sort the
points. On each open
piece $I=(l,r)$, the assigned plan, the Voice type, the crossing date, the filing date, and the
filing status are constant. Each payoff formula on that piece extends continuously to its closure
from the relevant side. The partition is finite, and its endpoints are null by (R-2). It therefore
suffices to bound every branch closure.

Write
\[
b=b^*(s),
\qquad
v_h=\mu_v+\beta(s-\mu_v),
\qquad
C=C_0\exp\!\left[-\frac{\chi(s-\mu_v)}{\sigma_s}\right].
\]
For Voice type $m$, let $e_m$ be the fixed pass's expected control-node entry probability,
$q_m$ its expected product of entry and control-node price, and $P_{dm}$ its expected pooled
execution price at date $d$. Set
\[
A_m=\frac1m\sum_{d=0}^{m-1}P_{dm}.
\]
Direct substitution into the blockholder's payoff gives the three branch formulas
\begin{align}
\label{eq:regret-branches-pooled}
U_{\mathsf{E}}&=b_0P_{00},
\nonumber\\
U_{\mathsf{H}}&=b_0\bigl((1-e_0)v_h+q_0+e_0m_0\bigr),
\\
Y_m&=(1-e_m)(v_h+\Delta_V)+q_m+e_m m_1,
\nonumber\\
U_{\mathsf{V}}&=bY_m-A_m(b-b_0)-C
\quad\text{on a pooled Voice branch}.
\nonumber
\end{align}
The Exit payoff is constant and the Hold payoff is affine. On a flagged Voice branch with filing
date $f$, put
\[
r_f=\min\{1,(f+1)/m\},
\qquad
B^F=b_0+r_f(b-b_0),
\qquad
A_{mf}=\frac1m\sum_{d=0}^{\min\{f,m-1\}}P_{dm}.
\]
At engagement posterior one, the flagged price satisfies
\[
P_F=(1-p_F)(v_h+\Delta_V)+p_F(P_F+m_1).
\]
This fixed point makes expected terminal value equal to $P_F$. The flagged Voice payoff therefore
reduces to
\begin{equation}
\label{eq:regret-branch-flagged}
U_{\mathsf{V}}=B^FP_F-A_{mf}(b-b_0)-C.
\end{equation}
Thus every payoff is given by a closed form on each branch, apart from the unique scalar price root
in \eqref{eq:regret-branch-flagged}, whose uniqueness follows from Lemma~\ref{lem:pricing-root}.

It remains to cover the space between mesh points. The target and engagement cost satisfy
\[
b'(s)=\frac{\bar b-b_0}{2\sigma_s}
\left(1+\left(\frac{s-\mu_v}{\sigma_s}\right)^2\right)^{-3/2},
\qquad
C'(s)=-\frac{\chi}{\sigma_s}C(s).
\]
Let $b'_{\max,I}$ be the maximum of $b'$ on the closed piece, attained at the point nearest
$\mu_v$. Differentiating the pooled Voice branch gives
\[
U_{\mathsf{V}}'=b'(Y_m-A_m)+b(1-e_m)\beta+\frac{\chi}{\sigma_s}C.
\]
Since $Y_m-A_m$ is affine, a valid Lipschitz constant on $I$ is
\begin{equation}
\label{eq:regret-LP}
L_{\mathsf{V}}^{P}
=b'_{\max,I}\max_{s\in\{l,r\}}|Y_m(s)-A_m|
+\bar b(1-e_m)\beta+\frac{\chi}{\sigma_s}C(l).
\end{equation}
The other two constants are $L_{\mathsf{E}}=0$ and
$L_{\mathsf{H}}=b_0(1-e_0)\beta$.

For the flagged branch, put $V=v_h+\Delta_V$ and
\[
a(P)=1-p(P)=\Phi\!\left(\frac{P+K+m_1-\bar S}{\sigma_\xi}\right),
\qquad a=a(P_F).
\]
Implicit differentiation of the price equation gives
\begin{equation}
\label{eq:regret-price-derivative}
\frac{dP_F}{ds}
=\beta\frac{a^2}{a^2+(m_1/\sigma_\xi)\phi_{\!N}(u)},
\qquad
0<\frac{dP_F}{ds}\le\beta,
\end{equation}
where $u=(P_F+K+m_1-\bar S)/\sigma_\xi$ and $\phi_{\!N}$ is the standard normal density. The root
also satisfies
\[
V\le P_F\le V+m_1\frac{1-a(V)}{a(V)}.
\]
Since $V$ rises in $s$ while $(1-a(V))/a(V)$ falls, the endpoint mixture
\[
P_{\mathrm{abs},I}
=\max\left\{|V(l)|,
\left|V(r)+m_1\frac{1-a(V(l))}{a(V(l))}\right|\right\}
\]
bounds $|P_F|$ on $I$. Differentiating \eqref{eq:regret-branch-flagged} then gives the valid constant
\begin{equation}
\label{eq:regret-LF}
L_{\mathsf{V}}^{F}
=b'_{\max,I}\bigl(r_fP_{\mathrm{abs},I}+|A_{mf}|\bigr)
+\bar b\beta+\frac{\chi}{\sigma_s}C(l).
\end{equation}

On each piece take $L_{\mathsf{V}}=L_{\mathsf{V}}^P$ or
$L_{\mathsf{V}}=L_{\mathsf{V}}^F$ according to its filing status.
Let $a_I$ be the benchmark plan on a piece and let
$G_{j,a_I}=U_j-U_{a_I}$ for either alternative $j$. The preceding bounds make
$G_{j,a_I}$ Lipschitz on the branch closure with constant
$L_{j,a_I}=L_j+L_{a_I}$. Divide $[l,r]$ into
\[
n_I=\left\lceil\frac{r-l}{10^{-5}}\right\rceil
\]
equal segments, and call the endpoint mesh $\mathcal G_I$. Every point of the piece lies within
$\delta_I=(r-l)/(2n_I)$ of a mesh point. Hence
\begin{equation}
\label{eq:regret-cover}
\sup_{s\in I}G_{j,a_I}(s)
\le \max_{s\in\mathcal G_I}G_{j,a_I}(s)+L_{j,a_I}\delta_I.
\end{equation}
The certificate is rigorous in exact arithmetic. It adds $5\times10^{-12}$ to the right side for
floating-point evaluation and the reported price-root residual. Define $\overline R$ as the maximum
of zero and these adjusted bounds over every piece and
both alternative plans. Equation~\eqref{eq:regret-cover} then bounds $R(s)$ on every open piece.
The omitted endpoints form a null set, so it also bounds the essential supremum on $I_s$. This is
the first inequality in \eqref{eq:regret-bound}. Dividing by the positive normaliser $N$ proves the
second.
\end{proof}

The bound establishes nothing about equilibrium existence, order size, speed, or a reacting
blockholder.


% proofs/06_lemmas.tex
% Two lemmas supporting the pricing root and timing claims.
% Fragment, to be \input by appendix.tex.

\section{Pricing root and timing lemmas}
\label{sec:pricing-timing-lemmas}

This section states and proves two lemmas the main text uses: the competitive pricing root is
unique and continuous, and at fixed policies a shorter window weakly lowers the stake at filing
on every path flagged under both windows.

\begin{lemma}[Uniqueness and continuity of the competitive pricing root]
\label{lem:pricing-root}
Assume the standing conditions \ref{p:space} and \ref{p:kernel}; the latter gives an entry
probability in $(0,1)$ that is continuous in the posterior and the price, and this lemma adds its
normal form and its monotonicity in the price. Let $m_0 \ge 0$, $\Delta_m > 0$, $\sigma_\xi > 0$,
$\Delta_V \ge 0$, and let the potential bidder's entry probability be
\begin{equation}
\label{eq:bidder-entry-prob}
p(P, \pi) \;=\; 1 - \Phi\!\left(\frac{P + K + m_0 + \pi\Delta_m - \bar S}{\sigma_\xi}\right),
\end{equation}
where $\Phi$ is the standard normal cumulative distribution function and $K, \bar S$ are finite constants. Then, for every $(\hat v, \pi) \in \mathbb{R} \times [0, 1]$, the competitive control-node pricing equation
\begin{equation}
\label{eq:pricing-fixed-point}
P \;=\; \bigl(1 - p(P, \pi)\bigr)\bigl(\hat v + \pi\Delta_V\bigr) \;+\; p(P, \pi)\bigl(P + m_0 + \pi\Delta_m\bigr)
\end{equation}
has a unique finite solution $P = \Pi(\hat v, \pi)$. The solution map $\Pi: \mathbb{R} \times [0, 1] \to \mathbb{R}$ is continuous, strictly increasing in $\hat v$, and satisfies the linear-growth bound
\begin{equation}
\label{eq:pricing-growth-bound}
|\Pi(\hat v, \pi)| \;\le\; |\hat v| + C_\Pi
\end{equation}
for the finite constant $C_\Pi := \max_{\pi \in [0, 1]} |\Pi(0, \pi)|$.
\end{lemma}

\begin{proof}
Set $\hat V = \hat v + \pi\Delta_V$ and $\widetilde m = m_0 + \pi\Delta_m$. Since $m_0 \ge 0$, $\Delta_m > 0$, and $\pi \in [0, 1]$, we have $\widetilde m \ge m_0 \ge 0$. Because $\Phi$ is strictly between $0$ and $1$ everywhere on the real line, $p(P, \pi) \in (0, 1)$ for every finite $P$. Rearranging \eqref{eq:pricing-fixed-point} by collecting terms in $P$ gives
\[
\bigl(1 - p(P, \pi)\bigr) P \;=\; \bigl(1 - p(P, \pi)\bigr)\hat V \;+\; p(P, \pi)\,\widetilde m,
\]
or equivalently, dividing by $1 - p(P, \pi) > 0$,
\begin{equation}
\label{eq:pricing-root-residual}
R(P; \hat V, \pi) \;:=\; P \;-\; \hat V \;-\; \widetilde m\,\frac{p(P, \pi)}{1 - p(P, \pi)} \;=\; 0.
\end{equation}
The entry probability $p(P, \pi)$ in \eqref{eq:bidder-entry-prob} is strictly decreasing and continuously differentiable in $P$, with $\lim_{P \to -\infty} p(P, \pi) = 1$ and $\lim_{P \to +\infty} p(P, \pi) = 0$. The odds ratio $q \mapsto q / (1 - q)$ is strictly increasing on $(0, 1)$. Therefore, $P \mapsto p(P, \pi) / (1 - p(P, \pi))$ is strictly decreasing in $P$. Since $\widetilde m \ge 0$, the function
\[
P \;\longmapsto\; -\,\widetilde m\,\frac{p(P, \pi)}{1 - p(P, \pi)}
\]
is weakly increasing in $P$. Consequently, $R(\cdot; \hat V, \pi)$ is strictly increasing on $\mathbb{R}$.

Moreover, as $P \to -\infty$, $p(P, \pi) \to 1$ and $p/(1-p) \to +\infty$, so $R(P; \hat V, \pi) \to -\infty$. As $P \to +\infty$, $p(P, \pi) \to 0$ and $p/(1-p) \to 0$, so $R(P; \hat V, \pi) \to +\infty$. By the intermediate value theorem, a solution to $R(P; \hat V, \pi) = 0$ exists. By strict monotonicity of $R$ in $P$, the solution is unique.

To establish continuity, differentiate $R$ with respect to $P$:
\[
\frac{\partial R}{\partial P} \;=\; 1 \;-\; \widetilde m\,\frac{\partial}{\partial P}\left(\frac{p}{1-p}\right) \;\ge\; 1 \;>\; 0,
\]
since $\widetilde m \ge 0$ and $\partial_P (p/(1-p)) < 0$. By the implicit function theorem, the unique root $\Pi(\hat v, \pi)$ is continuously differentiable in $(\hat V, \pi)$, and therefore continuous in $(\hat v, \pi)$ on $\mathbb{R} \times [0, 1]$. Furthermore,
\[
\frac{\partial \Pi}{\partial \hat V} \;=\; \frac{1}{\partial R / \partial P} \;\in\; (0, 1],
\]
so $\Pi$ is strictly increasing in $\hat v$. Write $\Pi^\ast(\hat V, \pi)$ for the root as a function of $\hat V$, so that $\Pi(\hat v, \pi) = \Pi^\ast(\hat v + \pi\Delta_V, \pi)$. For fixed $\pi$, the derivative bound above makes $\Pi^\ast$ $1$-Lipschitz in $\hat V$, and the two arguments $\hat v + \pi\Delta_V$ and $\pi\Delta_V$ differ by $\hat v$, so $|\Pi(\hat v, \pi) - \Pi(0, \pi)| \le |\hat v|$. Since $[0, 1]$ is compact and $\pi \mapsto \Pi(0, \pi)$ is continuous, $C_\Pi := \max_{\pi \in [0, 1]} |\Pi(0, \pi)| < \infty$. The bound \eqref{eq:pricing-growth-bound} follows.
\end{proof}

\begin{lemma}[Window monotonicity of the stake at filing under fixed policies]
\label{lem:bf-monotone}
Assume the standing conditions \ref{p:menu}, \ref{p:paths}, \ref{p:clock}, \ref{p:nofeedback}, and \ref{p:fixed}. Let the plan menu, the cutoff vector, and the execution paths $(B_j(s, \cdot))_{j \in \mathcal J}$ be fixed. Compare two window margins $T' < T$ at a common threshold $\tau$. Write $c(s;\tau)$ for the crossing date of a Voice type, $f(s;T_0) = c(s;\tau) + T_0$ for the filing date under window $T_0$, and $B^F(s;T_0) = B_{\mathrm{Voice}}(s, f(s;T_0))$ for the stake at filing, defined on the flagged cell $\{f(s;T_0) \le H\}$ and nowhere else. Then, for every signal realization $s$:
\begin{enumerate}[label=(\alph*),leftmargin=2.4em,itemsep=3pt]
\item On Exit and Hold plans no filing lands under either window margin, and the stake path is the same under both.
\item On Voice plans, the crossing date $c(s;\tau)$ is the same under both window margins, and the flagged cells are nested: $\{f(s;T) \le H\} \subseteq \{f(s;T') \le H\}$. On the difference $\{c(s;\tau) + T' \le H < c(s;\tau) + T\}$ the filing lands under $T'$ and does not land under $T$, so $B^F(s;T')$ is defined and $B^F(s;T)$ is not.
\item On Voice plans with $c(s;\tau) + T \le H$, the paths flagged under both windows, the filing lands weakly earlier under $T'$, $f(s;T') < f(s;T) \le H$, and
\begin{equation}
\label{eq:bf-monotone}
B^F(s; T') \;\le\; B^F(s; T).
\end{equation}
\end{enumerate}
In particular, shortening the window weakly lowers the stake at filing on every path flagged under both windows, and weakly enlarges the set of paths on which a filing lands.
\end{lemma}

\begin{proof}
Under fixed policies (\ref{p:fixed}) the plan selection and the stake paths are identical under $T$ and $T'$. By \ref{p:clock} only Voice plans cross $\tau$, so on Exit and Hold plans no filing lands under either window and the stake path is common; this is part~(a).

For a Voice type, \ref{p:clock} and \ref{p:nofeedback} make the crossing date $c(s;\tau) = \min\{d \le H : B_{\mathrm{Voice}}(s,d) \ge \tau\}$ (or $+\infty$) a function of the signal and $\tau$ alone, so it is the same under both windows. By \ref{p:clock} the filing lands exactly at $f(s;T_0) = c(s;\tau) + T_0$, and the flagged cell is the event $f(s;T_0) \le H$. Since $T' < T$, $f(s;T') < f(s;T)$ whenever $c(s;\tau) < \infty$, so $f(s;T) \le H$ implies $f(s;T') \le H$: the flagged cells are nested. On $\{c + T' \le H < c + T\}$ the filing lands under $T'$ and not under $T$, which is the last clause of part~(b).

For part~(c), let $c(s;\tau) + T \le H$. Then both filing dates lie in the calendar, $f(s;T') < f(s;T) \le H$, and both stakes at filing are defined. By \ref{p:paths} the Voice stake path $d \mapsto B_{\mathrm{Voice}}(s,d)$ is weakly increasing in the calendar date, so $B^F(s;T') = B_{\mathrm{Voice}}(s, f(s;T')) \le B_{\mathrm{Voice}}(s, f(s;T)) = B^F(s;T)$, which is \eqref{eq:bf-monotone}.
\end{proof}


\end{document}
